\documentclass[11pt]{amsart}
\usepackage[utf8]{inputenc}
\usepackage[T1]{fontenc}
\usepackage{lmodern}
\usepackage[margin=1.05in]{geometry}
\usepackage{amsmath,amssymb,amsthm,mathtools}
\usepackage{booktabs}
\usepackage{array}
\usepackage{enumitem}
\usepackage{listings}
\usepackage{xcolor}
\usepackage[hidelinks]{hyperref}
\usepackage{url}
\usepackage{microtype}
\makeatletter\g@addto@macro\UrlBreaks{\do\_\do\-}\makeatother
\emergencystretch=2em

\newtheorem{theorem}{Theorem}[section]
\newtheorem{proposition}[theorem]{Proposition}
\newtheorem{lemma}[theorem]{Lemma}
\newtheorem{corollary}[theorem]{Corollary}
\newtheorem{conjecture}[theorem]{Conjecture}
\theoremstyle{definition}
\newtheorem{definition}[theorem]{Definition}
\newtheorem{remark}[theorem]{Remark}
\numberwithin{equation}{section}

\newcommand{\Z}{\mathbb{Z}}
\newcommand{\Q}{\mathbb{Q}}
\newcommand{\N}{\mathbb{N}}
\newcommand{\R}{\mathcal{R}}
\newcommand{\Mel}{\mathrm{Mel}}
\newcommand{\qbin}[2]{\genfrac{[}{]}{0pt}{}{#1}{#2}}
\newcommand{\poch}[2]{(#1;q)_{#2}}
\newcommand{\repo}{\nolinkurl{2026-11-number-theory-q-series/}}
\newcommand{\fpath}[1]{\nolinkurl{#1}}

\lstdefinelanguage{lean}{
  morekeywords={theorem,def,noncomputable,variable,fun,if,then,else,by,let,in},
  sensitive=true, morecomment=[l]{--},
}
\lstset{language=lean,basicstyle=\ttfamily\footnotesize,breaklines=true,columns=fullflexible,keepspaces=true,
  keywordstyle=\bfseries,commentstyle=\itshape\color{black!60},
  literate={ℚ}{{$\mathbb{Q}$}}1 {ℕ}{{$\mathbb{N}$}}1 {⟦}{{$[\![$}}1 {⟧}{{$]\!]$}}1
   {∑}{{$\sum$}}1 {∏}{{$\prod$}}1 {∈}{{$\in$}}1 {∉}{{$\notin$}}1 {→}{{$\to$}}1
   {≤}{{$\le$}}1 {⁻¹}{{$^{-1}$}}2 {∀}{{$\forall$}}1 {∧}{{$\wedge$}}1 {∨}{{$\vee$}}1
   {⟨}{{$\langle$}}1 {⟩}{{$\rangle$}}1 {·}{{$\cdot$}}1 {−}{{$-$}}1 {∞}{{$\infty$}}1 {×}{{$\times$}}1}

\title[The four Andrews--Gordon-type identities of KRTW]{The four Andrews--Gordon-type identities of Kanade,
Russell, Tsuchioka and Warnaar: proofs for every level, with the first formalized in Lean}
\author{Raghav Sharma}
\email{raghavsharma1729@gmail.com}
\thanks{\textbf{AI disclosure.} The results, proofs, computer code and Lean formalizations in this
paper were produced by AI (Anthropic's Claude models) working under the direction of the author.
Correctness was checked by independent AI referees (separate model instances instructed to find errors,
using their own code; two independent reads of every proof, see §\ref{ssec:reads}) and, for the formalized
statements (identity~(i) only), by the Lean kernel. No human expert has refereed this work yet. Code,
certificates and the Lean project:
\url{https://github.com/raghavs1729/Research-Math}, folder \texttt{2026-11-number-theory-q-series/}.}
\date{27 November 2026}

\subjclass[2020]{Primary 11P84; Secondary 05A17, 05A30, 33D15, 05E05, 68V20}
\keywords{Rogers--Ramanujan identities, Andrews--Gordon identities, Bailey's lemma, multisums,
Hall--Littlewood polynomials, Capparelli--Meurman--Primc--Primc conjectures, Lean, formal verification}

\begin{document}

\begin{abstract}
Kanade, Russell, Tsuchioka and Warnaar (Selecta Math.\ 2026) conjectured four Andrews--Gordon-type
identities for $2k$-fold sums in the bases $q$ and $q^2$; for $k\le2$ they follow from known results. We prove
all four for every $k$. For the first, the $2k$-fold sum is reduced to a known $k$-fold sum through a finite
identity with a boundary parameter, proved from the $q$-binomial theorem; this identity, in the printed
product form and for every $k$, is formalized in Lean~4 with Mathlib. For the other three we introduce a
parameter $x$ and write the $2k$-fold sum as a matrix element of a transfer operator on functions of two
variables. We diagonalise this operator without computing its eigenfunctions: two commuting difference
operators give one-step recursions for the values of the eigenfunctions at infinity, the left and right
eigenfunctions of each mode pair at a single point, and the resulting spectral expansion is a product of
two Bailey-type series. This proves the factorisation $K_k(x)=B_{k+1}(x)J_k(x)$ conjectured in the first version of
this paper, and hence the second and third identities. For the fourth, the inserted weight is a
commutator with the transfer operator up to a factor, so its insertions telescope, and the same spectral
expansion applies. With theorems of Griffin, Ono and Warnaar this proves KRTW's Hall--Littlewood conjecture
at $z=1$ and at $z=q$ for every $k$. The results, proofs, code and Lean formalization were produced by AI
(Anthropic's Claude models) directed by the author; they were checked by independent AI referees and, for
the first identity, by the Lean kernel. No human expert has refereed this work yet.
\end{abstract}

\maketitle


\section{Introduction}\label{sec:intro}

\subsection{The conjectures}
Throughout, $(a;q)_n=\prod_{j=0}^{n-1}(1-aq^j)$, $(q)_n=(q;q)_n$,
$(a_1,\dots,a_s;q)_\infty=\prod_i\prod_{j\ge0}(1-a_iq^j)$, and $1/(q)_n=1/(q^2;q^2)_n:=0$ for $n<0$. Following~\cite{GOW}
we write $\theta(a;p)=(a,p/a;p)_\infty$ and $\theta(a_1,\dots,a_s;p)=\theta(a_1;p)\cdots\theta(a_s;p)$; the
Jacobi triple product reads $\sum_{n\in\Z}(-1)^nz^np^{n(n-1)/2}=(p;p)_\infty\theta(z;p)$.

In their study of the conjectures of Capparelli, Meurman, Primc and Primc on coloured partitions,
Kanade, Russell, Tsuchioka and Warnaar~\cite{KRTW} (henceforth KRTW) introduced the $2k$-fold
weight
\begin{equation}\label{eq:W}
W_k(r,s)=\prod_{i=1}^{k}\frac{q^{(r_i+s_i)^2+s_i^2}}{(q;q)_{r_i-r_{i+1}}\,(q^2;q^2)_{s_i-s_{i-1}}},
\qquad r_{k+1}=s_0:=0,
\end{equation}
summed over all $r=(r_1,\dots,r_k)\in\N^k$, $s=(s_1,\dots,s_k)\in\N^k$; by the convention
$1/(q)_n=0$ for $n<0$ the sum runs over $r_1\ge\dots\ge r_k\ge0$ and $0\le s_1\le\dots\le s_k$. After
their Conjecture~4.9 they display four ``conjectural Andrews--Gordon-type identities''
\cite[§4, display after Conjecture~4.9]{KRTW}:
\begin{align}
\textstyle\sum q^{\sum_i s_i}\,W_k(r,s)&=\frac{(q,q^{k+1},q^{k+2};q^{k+2})_\infty}{(q;q^2)_\infty(q;q)_\infty},
\tag{i}\label{eq:i}\\
\textstyle\sum q^{\sum_i(r_i+2s_i)}\,W_k(r,s)&=\frac{(q,q^2,q^{2k+2},q^{2k+3},q^{2k+4},q^{2k+4};q^{2k+4})_\infty}{(q^2;q^2)_\infty(q;q)_\infty},
\tag{ii}\label{eq:ii}\\
\textstyle\sum W_k(r,s)&=\frac{(q^{k+1},q^{k+2},q^{k+2},q^{k+3},q^{2k+4},q^{2k+4};q^{2k+4})_\infty}{(q^2;q^2)_\infty(q;q)_\infty},
\tag{iii}\label{eq:iii}\\
\textstyle\sum \Omega\,W_k(r,s)&=\frac{(q^{k},q^{k+1},q^{k+3},q^{k+4},q^{2k+4},q^{2k+4};q^{2k+4})_\infty}{(q^2;q^2)_\infty(q;q)_\infty},
\tag{iv}\label{eq:iv}
\end{align}
where $\Omega=\sum_{i<k}q^{r_i+2s_i}(1-q^{2s_{i+1}-2s_i})+q^{r_k+2s_k}$. Identity~\eqref{eq:i} is also
displayed in the introduction of~\cite{KRTW} among ``a number of new conjectures of Rogers--Ramanujan
type''. Its product side is the principally specialised character of the level-$k$ standard module
$L(k\Lambda_0)$ of $A_1^{(1)}$. KRTW derive the four identities from their Conjecture~4.9 (a
multisum formula for the generating functions of the $D_3^{(2)}$-type coloured partitions), their
identity (4.7), and the product formulas of Conjectures~4.1 and~4.2. Identity~\eqref{eq:i} is,
moreover, the case $z=1$ of their Hall--Littlewood conjecture~\cite[Conjecture~2.2]{KRTW}:
\begin{equation}\label{eq:conj22}
\sum_{\lambda:\ \lambda_1\le k}(zq)^{|\lambda|}P_{2\lambda}(1,q,q^2,\dots;q^2)
=\sum_{r,s}\prod_{i=1}^k\frac{z^{r_i+s_i}q^{(r_i+s_i)^2+s_i^2+s_i}}{(q;q)_{r_i-r_{i+1}}(q^2;q^2)_{s_i-s_{i-1}}},
\end{equation}
where $P_\lambda(\,\cdot\,;t)$ is the Hall--Littlewood polynomial and the sum on the left is over
partitions $\lambda$ with largest part at most $k$.

\subsection{What was known}\label{ssec:known}
For $k=1$, identities \eqref{eq:i}, \eqref{eq:ii}, \eqref{eq:iv} follow from results of KRTW
(Conjecture~2.2 for $k\le1$, Proposition~4.8) and classical level-one product formulas; for
\eqref{eq:iii} at $k=1$ the sum is, term by term, a Bressoud-type double sum (see
Remark~\ref{rem:k1}), but KRTW do not discuss this case separately. For $k=2$, identity~\eqref{eq:i}
follows from KRTW's Theorem~4.10 (the multisum formulas for $k=2$, proved by functional equations), their
identity~(4.7) and the theorem of Meurman and Primc~\cite{MP} (which, as KRTW state, establishes their
Conjecture~4.1 for $n=1$); identities \eqref{eq:ii}--\eqref{eq:iv} for $k=2$ follow from KRTW's
Theorem~4.10 together with Bartholdi's recent proof of the three-row, level-two case of Primc's
conjecture~\cite{Bartholdi} (this deduction is ours; it is not stated in either paper). For $k\ge3$ we
found no proof of any of the four identities in the literature (see §\ref{ssec:novelty}).

\subsection{Main results}
Our first result settles identity~\eqref{eq:i}.

\begin{theorem}\label{thm:main}
Identity~\eqref{eq:i} holds for every integer $k\ge0$:
\[
\sum_{\substack{r_1\ge\dots\ge r_k\ge0\\ 0\le s_1\le\dots\le s_k}}\ \prod_{i=1}^k
\frac{q^{(r_i+s_i)^2+s_i^2+s_i}}{(q;q)_{r_i-r_{i+1}}(q^2;q^2)_{s_i-s_{i-1}}}
=\frac{(q,q^{k+1},q^{k+2};q^{k+2})_\infty}{(q;q^2)_\infty(q;q)_\infty}.
\]
\end{theorem}

The proof has two halves. The first, which is the new part, is the following reduction (with an
auxiliary grading~$x$) of the $2k$-fold sum to a $k$-fold sum.

\begin{theorem}\label{thm:x}
For every $k\ge0$, in $\Z[x][[q]]$,
\begin{equation}\label{eq:xref}
\sum_{\substack{r_1\ge\dots\ge r_k\ge0\\0\le s_1\le\dots\le s_k}} x^{\sum_i(r_i+2s_i)}\prod_{i=1}^k
\frac{q^{(r_i+s_i)^2+s_i^2+s_i}}{(q)_{r_i-r_{i+1}}(q^2;q^2)_{s_i-s_{i-1}}}
=\sum_{M_1\ge\dots\ge M_k\ge0}\frac{x^{M_1+\dots+M_k}\,q^{\sum_iM_i(M_i+1)/2}}{\prod_{i=1}^k(q)_{M_i-M_{i+1}}}
\end{equation}
($M_{k+1}:=0$), and both sides equal
\begin{equation}\label{eq:melclosed}
\frac{(-xq;q)_\infty}{(x^2q;q)_\infty}\sum_{n\ge0}(-1)^nx^{(k+1)n}q^{(k+1)n(n+1)/2+n(n-1)/2}\,
\frac{(x^2q;q)_n\,(1-x^2q^{2n+1})}{(q;q)_n}.
\end{equation}
\end{theorem}

The first equality~\eqref{eq:xref} is deduced from a finite identity with a boundary parameter~$R$
(Theorem~\ref{thm:R}), whose key input is a polynomial identity (Lemma~\ref{lem:key}) proved from the
$q$-binomial theorem. The second equality is a Bailey-chain computation; at $x=1$ it is the known
evaluation of the $k$-fold sum (see §\ref{ssec:novelty}). Theorem~\ref{thm:main} is the case $x=1$,
combined with the Jacobi triple product. By a theorem of Griffin, Ono and Warnaar~\cite{GOW} we obtain:

\begin{corollary}\label{cor:z1}
KRTW's Conjecture~2.2, identity~\eqref{eq:conj22}, holds at $z=1$ for every $k\ge0$.
\end{corollary}

The second group of results concerns \eqref{eq:ii}--\eqref{eq:iv}. Write $Q=q^2$ and put
\begin{align*}
K_k(x)&:=\sum_{r,s}x^{\sum_i(r_i+2s_i)}\,W_k(r,s),\\
B_m(x)&:=\sum_{N_1\ge\dots\ge N_m\ge0}\frac{x^{N_1+\dots+N_m}q^{N_1^2+\dots+N_m^2}}{(q)_{N_1-N_2}\cdots(q)_{N_{m-1}-N_m}(Q;Q)_{N_m}},\\
C_k(x)&:=\sum_{0\le s_1\le\dots\le s_k}\frac{x^{s_1+\dots+s_k}q^{s_1^2+\dots+s_k^2}}{\prod_{i=1}^k(Q;Q)_{s_i-s_{i-1}}},
\qquad J_k(x):=\frac{C_k(x)}{C_1(x)},\quad C_1(x)=(-xq;Q)_\infty .
\end{align*}
Thus \eqref{eq:iii} and \eqref{eq:ii} are the evaluations of $K_k(1)$ and $K_k(q)$. The first version of this
paper conjectured the following factorisation, tested it for $k\le7$, and proved it for $k\le3$.

\begin{theorem}[$F_k$]\label{thm:F}
For every $k\ge1$, in $\Z[x][[q]]$,
\[K_k(x)=B_{k+1}(x)\,J_k(x),\qquad\text{equivalently}\qquad K_k(x)\,C_1(x)=B_{k+1}(x)\,C_k(x).\]
\end{theorem}

\begin{theorem}\label{thm:k3}
Identities \eqref{eq:ii} and \eqref{eq:iii} hold for every $k\ge1$.
\end{theorem}

For the fourth identity let $K^\Omega_k(x)$ be $K_k(x)$ with the factor $\Omega$ inserted in the summand, and let
$B^{(k)}_{k+1}(x)$ be $B_{k+1}(x)$ with the additional factor $q^{N_k+N_{k+1}}$ in the summand. Thus \eqref{eq:iv} is the
evaluation of $K^\Omega_k(1)$.

\begin{theorem}\label{thm:iv}
For every $k\ge1$, $K^\Omega_k(x)=B^{(k)}_{k+1}(x)\,J_k(x)$ in $\Z[x][[q]]$, and identity~\eqref{eq:iv} holds.
\end{theorem}

The proofs are in §\ref{sec:k3}. There $K_k(x)$ is written as a matrix element $\omega(T^{k-1}u)$ of a transfer
operator $T$ acting on functions of two variables $(R,S)$. $T$ is triangular in a Bailey-type basis, with
eigenvalues $x^{N+M}q^{N^2+M^2}$ ($N\ge M\ge0$), but its eigenfunctions are not hypergeometric terms, and we do
not compute them. Two difference operators $H'$, $G$ that commute with $T$ give, on the lowest slice and read at
$R=\infty$, a one-step recursion for the value at infinity of each eigenfunction; the same works for the
transposed chain, whose eigenfunctions give the left eigenvectors; the left and right eigenvectors of a mode
meet in a single point; and completeness follows from $x$-adic valuations. The resulting spectral expansion of
$K_k$ is the double-sum form of $B_{k+1}J_k$ given by Bailey's lemma. For \eqref{eq:iv}, the insertion of $\Omega$
at a level is the operator $O=q^{R+2S}\cdot(S\mapsto S+1)$, and $Th-hT=-xOT$ for a rank-one operator~$h$ that already
occurs in $H'$; so the $k$ insertions telescope to a boundary term, whose spectral weights factor again.
Theorem~\ref{thm:k3} and a theorem of Griffin, Ono and Warnaar~\cite[Theorem~1.3]{GOW} give:

\begin{corollary}\label{cor:zq}
Identity~\eqref{eq:ii} for a given $k\ge1$ is equivalent to KRTW's Conjecture~2.2 at $z=q$ for that $k$. Hence
Conjecture~2.2 holds at $z=q$ for every $k\ge1$.
\end{corollary}

We also recall closed Bailey-type formulas for $B_{k+1}(x)$, $J_k(x)$ and their product (Theorem~\ref{thm:B}); the
product is the $C_2$ Rogers--Selberg series of Bartlett and Warnaar~\cite{BW} at the point $(\sqrt x,-\sqrt x)$,
so that, given their $C_2$ Rogers--Selberg identity, Theorem~\ref{thm:F} is a Hall--Littlewood identity $(HL_k)$
of the same shape as~\eqref{eq:conj22} along the grading~$x$ (Corollary~\ref{cor:HL}).

\subsection{Formalization}
Theorem~\ref{thm:main}, in the product form printed by KRTW and for every $k$, and the first equality
of Theorem~\ref{thm:x} for an arbitrary power series~$x$, are proved in Lean~4 with
Mathlib~\cite{Lean4,Mathlib}; the main theorem is \fpath{KRTW.krtw_identity}. The formalization uses no
\texttt{sorry}, no added axiom and no \texttt{native\_decide}; \verb|#print axioms| reports only
\texttt{propext}, \texttt{Classical.choice}, \texttt{Quot.sound}. A reader can therefore check
Theorem~\ref{thm:main} without trusting this paper: read the eleven short definitions of the statement
file against the display of~\cite{KRTW}, run \texttt{lake build}, and run the two supplied checks that the
proved statements are the trusted ones and depend only on the standard axioms (§\ref{sec:lean}). Only
identity~\eqref{eq:i} is formalized: Theorems~\ref{thm:F}, \ref{thm:k3}, \ref{thm:iv}, Theorem~\ref{thm:B} and the
corollaries are not; for them the evidence is the proofs, the independent referee reads (§\ref{ssec:reads}) and
the certified computations (§\ref{sec:num}).

\subsection{What is new, what was known, what is open}\label{ssec:novelty}
\emph{Known.} The evaluation of the reduced $k$-fold sum,
\[
\sum_{M_1\ge\dots\ge M_k\ge0}\frac{q^{\sum_iM_i(M_i+1)/2}}{\prod_i(q)_{M_i-M_{i+1}}}
=\frac{(q,q^{k+1},q^{k+2};q^{k+2})_\infty}{(q;q^2)_\infty(q;q)_\infty},
\]
is in the literature for every $k$: Chern, Li, Stanton, Xue and Yee~\cite[Theorem~3.5]{CLSXY};
Kanade and Lovejoy~\cite[Theorem~1.1 with $a=0$]{KL}, by Bailey pairs; Warnaar~\cite[display after
Corollary~5, case $i=0$]{Warnaar2026}, where it is attributed to Andrews~\cite[Theorem~3]{Andrews2010},
Kim and Yee~\cite[Theorem~1.4]{KimYee}, \cite{CLSXY} and Kanade and Russell~\cite[Theorem~13]{KR};
and Dousse and Jang~\cite{DJ} (their Andrews--Kim--Yee-type theorem with $a=0$). For even $k$ it is
also the case $s=2$ of Melzer's identities~\cite[(2.6)]{Melzer}, proved by Bressoud, Ismail and
Stanton~\cite[Theorem~5.2]{BIS} (this identification was checked exponent by exponent by one of the
referees). Hence
Proposition~\ref{prop:bailey} below is a re-proof; we did not find its $x$-refined form
\eqref{eq:melclosed} stated in these papers, but it is the same Bailey chain, and we claim no novelty
for it. The Bailey-type formulas (B1)--(B5) of Theorem~\ref{thm:B} and the change-of-$a$ step in the proof of
Theorem~\ref{thm:iv} are short consequences of Bailey's lemma; at $x=1$ they are Bressoud's even-modulus
identities~\cite{Bressoud79,Bressoud80}; we do not claim them as new. The Bailey-basis facts of §\ref{ssec:basis}
are classical.

\emph{New.} The first equality~\eqref{eq:xref} of Theorem~\ref{thm:x}, with its finite form
(Theorem~\ref{thm:R} and Lemma~\ref{lem:key}); hence Theorem~\ref{thm:main} and Corollary~\ref{cor:z1} for $k\ge3$.
Theorem~\ref{thm:F} for $k\ge2$ (the first version proved $k=2,3$ by chains of Euler and Durfee moves);
Theorem~\ref{thm:k3} and Corollary~\ref{cor:zq} for $k\ge3$ (the first version proved $k=3$);
Theorem~\ref{thm:iv} for $k\ge3$, and its $x$-refinement for $k\ge2$; the method of §\ref{sec:k3} (commuting
difference operators, values at infinity, single-point biorthogonality, and the telescoping of the
$\Omega$-insertions); the Hall--Littlewood form $(HL_k)$ for every $k$, given (RS) (Corollary~\ref{cor:HL}). The Lean
formalization of~\eqref{eq:i}.

\emph{Open.} KRTW's Conjecture~2.2 for general $z$ and their Conjecture~4.9 for $k\ge3$; Primc's $D_3^{(2)}$
conjecture at level $k\ge3$; a formalization of \eqref{eq:ii}--\eqref{eq:iv}.

\emph{Searches.} The novelty statements rest on searches of 26 and 27 November 2026: arXiv listings and
metadata searches up to the submissions of 25 November 2026 (queries on Capparelli--Meurman--Primc,
Hall--Littlewood and Rogers--Ramanujan, Andrews--Gordon, Primc's conjecture, Melzer, Bailey, and every arXiv
paper by Kanade, Russell, Tsuchioka or Warnaar posted after~\cite{KRTW}); the OpenAlex and Semantic Scholar
lists of works citing~\cite{KRTW}; and the full texts of the most relevant recent papers, among
them~\cite{Bartholdi,Ding,Kursungoz,Russell,Warnaar2026,KL,DJ}. None contains a proof of
\eqref{eq:i}--\eqref{eq:iv} for $k\ge3$. No full-text search engine (MathSciNet, zbMATH, Google
Scholar) was available, so an unindexed earlier proof cannot be excluded.

\subsection{Verification status}\label{ssec:reads}
Each proof was read by two AI referees: separate model instances instructed to find errors, each writing its
own code. The second referee of each proof formed its verdict independently of the first report (for §\ref{sec:k3}
it was not shown that report or its code). The reports are in the repository, the first four in
\fpath{01-krtw-identities/}, the others in \fpath{01-krtw-identities/all-k/}.
\begin{center}\footnotesize
\begin{tabular}{@{}p{0.34\textwidth}p{0.42\textwidth}p{0.16\textwidth}@{}}
\toprule
Result & referee reads (reports) & further checks\\
\midrule
Theorems~\ref{thm:main}, \ref{thm:x}, \ref{thm:R}; Corollary~\ref{cor:z1} & 2: \fpath{REFEREE-KRTW-general.md},\newline \fpath{REFEREE-2-KRTW.md} & Lean (for~(i))\\
Theorem~\ref{thm:B} & 2: \fpath{REFEREE-KRTW-F-moves.md},\newline \fpath{REFEREE-2-KRTW-F.md} & certificates\\
Theorems~\ref{thm:F}, \ref{thm:k3} (§§\ref{ssec:transfer}--\ref{ssec:proofF}) & 2: \fpath{REFEREE-F-all-k.md},\newline \fpath{REFEREE-3-all-k.md} & certificates\\
Theorem~\ref{thm:iv} (§\ref{ssec:iv}) & 2: \fpath{REFEREE-iv-all-k.md},\newline \fpath{REFEREE-3-all-k.md} & certificates\\
\bottomrule
\end{tabular}
\end{center}
None of the reads found an error or a gap in the proofs of the results listed; the defects they raised (in side
remarks, the exposition and the literature discussion) were corrected. Corollaries~\ref{cor:z1}, \ref{cor:zq}, \ref{cor:y} and~\ref{cor:HL} rest in addition on the
cited theorems of Griffin, Ono and Warnaar and of Bartlett and Warnaar, which were not re-checked.

\subsection{A note on versions}
This is version~2 (27 November 2026). Version~1 (the same day) proved \eqref{eq:i} for every $k$, proved
\eqref{eq:ii}, \eqref{eq:iii} for $k\le3$ by exact chains of Euler and Durfee moves found by computer search, and
stated $(F_k)$ and its $\Omega$-analogue as conjectures. Version~2 replaces that part (old §7 and the appendix) by
§\ref{sec:k3}, which proves both for every $k$; §§\ref{sec:setting}--\ref{sec:cor} and the Lean section are
unchanged apart from cross-references. The move chains for $k\le3$ remain in the repository, with their certificates, as an independent
check.

\subsection{Organisation}
§\ref{sec:setting} fixes the formal setting. §\ref{sec:lemmas} proves the four lemmas,
§\ref{sec:R} the finite identity, §\ref{sec:x} Theorems~\ref{thm:x} and~\ref{thm:main},
§\ref{sec:cor} the corollaries. §\ref{sec:k3} treats \eqref{eq:ii}--\eqref{eq:iv}: the closed forms and the reduction to
identities in $x$ (§\ref{ssec:closed}), the transfer operator and its spectral theory
(§§\ref{ssec:transfer}--\ref{ssec:spec}), the proof of Theorem~\ref{thm:F} (§\ref{ssec:proofF}), identity~\eqref{eq:iv}
(§\ref{ssec:iv}) and the Hall--Littlewood form (§\ref{ssec:HL}). §\ref{sec:lean} describes the formalization,
§\ref{sec:num} the numerical checks.

\section{Formal setting}\label{sec:setting}

We work in $\R=\Z[x][[q]]$, the ring of formal power series in $q$ whose coefficients are polynomials
in $x$, with its $q$-adic topology, or in $\Q(q)[x]$ for finite identities. A family $(f_\iota)$ in
$\R$ is \emph{summable} if for every $n$ only finitely many $f_\iota$ are nonzero modulo $q^n$; summable
families may be regrouped and summed in any order. Every multiple series below is summable: its
general term is $x^aq^eu$ with $u\in\R$ and $e\to\infty$ along the index set (for instance
$e\ge\sum_i(r_i^2+s_i^2)$, or $e\ge\sum_iM_i(M_i+1)/2$). The substitutions $x\mapsto xq^c$ ($c\ge0$)
and $x\mapsto1$ are continuous ring homomorphisms $\R\to\R$, resp.\ $\R\to\Z[[q]]$. We use
$1/(q)_n,\ 1/(q^2;q^2)_n\in\Z[[q]]$, $(xq^c;q)_\infty\in\R$ for $c\ge1$, and the Gaussian
polynomial $\qbin nk=(q)_n/((q)_k(q)_{n-k})$ for $0\le k\le n$ ($:=0$ otherwise). For $N=r+2s$ we have
the elementary identity
\begin{equation}\label{eq:01}
(r+s)^2+s^2+s=\tfrac{N(N+1)}2+\tfrac{r(r-1)}2 .
\end{equation}

\section{Four lemmas}\label{sec:lemmas}

\begin{lemma}[$q$-binomial theorem; Euler]\label{lem:qbin}
For $n\ge0$ and an indeterminate $z$,
\[
(z;q)_n=\sum_{i=0}^n\qbin ni(-1)^iq^{i(i-1)/2}z^i,\qquad
\frac1{(z;q)_\infty}=\sum_{n\ge0}\frac{z^n}{(q)_n},\qquad
(z;q)_\infty=\sum_{n\ge0}\frac{(-1)^nq^{n(n-1)/2}z^n}{(q)_n}.
\]
\end{lemma}
These are \cite[(3.3.6), (2.2.5), (2.2.6)]{AndrewsBook}. The first is a polynomial identity in $z$
over $\Z[q]$, so $z$ may be specialised to any element of $\Z[q,q^{-1}][x]$.

\begin{lemma}\label{lem:L2}
For integers $R,L\ge0$,
\begin{equation}\label{eq:L2}
\sum_{\substack{t,b\ge0\\t+2b=L}}\frac{q^{t(t-1)/2}\qbin Rt}{(q^2;q^2)_b}
=\sum_{j=0}^{L}\frac{(-1)^jq^{Rj}}{(q)_j(q)_{L-j}} .
\end{equation}
\end{lemma}
\begin{proof}
By Lemma~\ref{lem:qbin} (with $z\mapsto-z$, and Euler's first series in base $q^2$) the left side is
the coefficient of $z^L$ in $(-z;q)_R/(z^2;q^2)_\infty$. Since
$(z^2;q^2)_\infty=(z;q)_\infty(-z;q)_\infty$ and $(-z;q)_R/(-z;q)_\infty=1/(-zq^R;q)_\infty$,
\[
\frac{(-z;q)_R}{(z^2;q^2)_\infty}=\frac1{(z;q)_\infty(-zq^R;q)_\infty}
=\sum_{l\ge0}\frac{z^l}{(q)_l}\sum_{j\ge0}\frac{(-q^Rz)^j}{(q)_j},
\]
whose coefficient of $z^L$ is the right side.
\end{proof}

\begin{lemma}\label{lem:L3}
For integers $R,d\ge0$ and an indeterminate $y$,
\begin{equation}\label{eq:L3}
\sum_{r=0}^{R}\qbin Rr y^rq^{r^2}(yq^{d+r+1};q)_{R-r}=\sum_{t=0}^{\min(R,d)}\qbin Rt\qbin dt(q)_t\,y^tq^{t^2}.
\end{equation}
\end{lemma}
\begin{proof}
Expand $(yq^{d+r+1};q)_{R-r}=\sum_i\qbin{R-r}i(-1)^iq^{i(i-1)/2}y^iq^{(d+r+1)i}$, put $t=r+i$ and use
$\qbin Rr\qbin{R-r}i=\qbin Rt\qbin ti$. With $r=t-i$ the exponent of $q$ is
$t^2+i(i-1)/2+i(d+1-t)$, so by Lemma~\ref{lem:qbin} the left side equals
\[\sum_t\qbin Rty^tq^{t^2}\sum_i\qbin ti(-1)^iq^{i(i-1)/2}(q^{d+1-t})^i=\sum_t\qbin Rty^tq^{t^2}(q^{d+1-t};q)_t.\] For $t>d$ the product $(q^{d+1-t};q)_t$ contains the factor $1-q^0=0$; for
$t\le d$ it equals $(q)_d/(q)_{d-t}=\qbin dt(q)_t$.
\end{proof}

\begin{lemma}[key identity]\label{lem:key}
For integers $R,M\ge0$, in $\Q(q)[x]$,
\begin{multline}\label{eq:star}
\sum_{r=0}^{R}\sum_{b=0}^{\lfloor M/2\rfloor}\qbin Rr
\frac{x^{r+2b}q^{(r+2b)(r+2b+1)/2+r(r-1)/2}\,(xq^{r+M+1};q)_{R-r}}{(q^2;q^2)_b\,(q)_{M-2b}}\\
=\sum_{L=0}^{M}\frac{x^Lq^{L(L+1)/2}\,(xq^{R+L+1};q)_{M-L}}{(q)_L\,(q)_{M-L}} .
\end{multline}
Consequently, in $\R$,
\begin{equation}\label{eq:star2}
\sum_{r=0}^{R}\sum_{b=0}^{\lfloor M/2\rfloor}\qbin Rr
\frac{x^{r+2b}q^{(r+2b)(r+2b+1)/2+r(r-1)/2}\,(xq^{r+M+1};q)_\infty}{(q^2;q^2)_b\,(q)_{M-2b}}
=\sum_{L=0}^{M}\frac{x^Lq^{L(L+1)/2}\,(xq^{R+L+1};q)_\infty}{(q)_L\,(q)_{M-L}},
\end{equation}
and
\begin{equation}\label{eq:starinf}
\sum_{r=0}^{R}\sum_{b\ge0}\qbin Rr\frac{x^{r+2b}q^{(r+2b)(r+2b+1)/2+r(r-1)/2}}{(q^2;q^2)_b}
=\sum_{L\ge0}\frac{x^Lq^{L(L+1)/2}(xq^{R+L+1};q)_\infty}{(q)_L} .
\end{equation}
\end{lemma}
\begin{proof}
We bring both sides of \eqref{eq:star} to the form
\begin{equation}\label{eq:23}
\sum_{L=0}^{M}\frac{x^Lq^{L(L+1)/2}}{(q)_{M-L}}\sum_{j=0}^{L}\frac{(-1)^jq^{Rj}}{(q)_j(q)_{L-j}} .
\end{equation}
\emph{Left side.} By \eqref{eq:01} the exponent is $r^2+2br+2b^2+b$, so the left side is
\[
\sum_b\frac{x^{2b}q^{2b^2+b}}{(q^2;q^2)_b(q)_{M-2b}}\sum_r\qbin Rr(xq^{2b})^rq^{r^2}\bigl(xq^{2b}\cdot q^{(M-2b)+r+1};q\bigr)_{R-r}.
\]
By Lemma~\ref{lem:L3} with $y=xq^{2b}$, $d=M-2b\ge0$, the inner sum is
$\sum_t\qbin Rt\qbin{M-2b}t(q)_tx^tq^{2bt+t^2}$, and $\qbin{M-2b}t(q)_t/(q)_{M-2b}=1/(q)_{M-2b-t}$. Hence
the left side equals
\[
\sum_{\substack{t,b\ge0\\t+2b\le M}}\frac{x^{t+2b}q^{t^2+2bt+2b^2+b}\qbin Rt}{(q^2;q^2)_b(q)_{M-2b-t}}
=\sum_{L=0}^{M}\frac{x^Lq^{L(L+1)/2}}{(q)_{M-L}}\sum_{t+2b=L}\frac{q^{t(t-1)/2}\qbin Rt}{(q^2;q^2)_b}
\]
(with $L=t+2b$ and \eqref{eq:01}), which is \eqref{eq:23} by Lemma~\ref{lem:L2}.
\emph{Right side.} Expanding $(xq^{R+L+1};q)_{M-L}$ by Lemma~\ref{lem:qbin} and using
$L(L+1)/2+i(i-1)/2+(L+1)i=(L+i)(L+i+1)/2$, the right side is
\[
\sum_{\substack{L,i\ge0\\L+i\le M}}\frac{(-1)^iq^{Ri}x^{L+i}q^{(L+i)(L+i+1)/2}}{(q)_L(q)_i(q)_{M-L-i}},
\]
which is \eqref{eq:23} after renaming $(L+i,i)\to(L,j)$. This proves \eqref{eq:star}. For
\eqref{eq:star2} multiply by $(xq^{R+M+1};q)_\infty$ and note that
$(xq^{r+M+1})_{R-r}(xq^{R+M+1})_\infty=(xq^{r+M+1})_\infty$ and
$(xq^{R+L+1})_{M-L}(xq^{R+M+1})_\infty=(xq^{R+L+1})_\infty$. For \eqref{eq:starinf}: by \eqref{eq:01}
and Lemma~\ref{lem:L2} the left side is
\[
\sum_Lx^Lq^{L(L+1)/2}\sum_{j+l=L}\frac{(-1)^jq^{Rj}}{(q)_j(q)_l}
=\sum_l\frac{x^lq^{l(l+1)/2}}{(q)_l}\sum_j\frac{(-1)^jq^{j(j-1)/2}(xq^{R+l+1})^j}{(q)_j},
\]
because $(j+l)(j+l+1)/2=l(l+1)/2+j(j-1)/2+j(l+1)$; the inner sum is $(xq^{R+l+1};q)_\infty$ by
Lemma~\ref{lem:qbin}.
\end{proof}

Lemma~\ref{lem:L2} for $R\to\infty$ is a classical self-fusion identity; its point here is that the
boundary $R$ enters the right side only through the factors $q^{Rj}$. Lemma~\ref{lem:L3} is a
terminating identity of $q$-Chu--Vandermonde type (for $d=0$ it is
$\sum_r\qbin Rry^rq^{r^2}/(yq)_r=1/(yq)_R$); we expect Lemmas~\ref{lem:L2} and~\ref{lem:L3} to be classical
in some form and claim novelty only for their combination in Lemma~\ref{lem:key} and its use below.

\section{The finite identity}\label{sec:R}

\begin{theorem}\label{thm:R}
For integers $j\ge1$ and $R\ge0$, in $\R$,
\begin{multline*}
A_j(R;x):=\sum_{\substack{R\ge r_1\ge\dots\ge r_j\ge0\\0\le s_1\le\dots\le s_j}}\frac{(q)_R}{(q)_{R-r_1}}
\prod_{i=1}^{j}\frac{x^{r_i+2s_i}q^{(r_i+s_i)^2+s_i^2+s_i}}{(q)_{r_i-r_{i+1}}(q^2;q^2)_{s_i-s_{i-1}}}\\
=D_j(R;x):=\sum_{M_1\ge\dots\ge M_j\ge0}\frac{x^{M_1+\dots+M_j}\,q^{\sum_iM_i(M_i+1)/2}\,(xq^{R+M_j+1};q)_\infty}{\prod_{i=1}^{j}(q)_{M_i-M_{i+1}}}
\end{multline*}
($r_{j+1}=s_0=M_{j+1}:=0$). Only the innermost variable $M_j$ feels the bound $R$.
\end{theorem}
\begin{proof}
Induction on $j$. For $j=1$, $(q)_R/((q)_{R-r}(q)_r)=\qbin Rr$ and \eqref{eq:01} show that $A_1(R;x)$ is the
left side of \eqref{eq:starinf} and $D_1(R;x)$ is its right side.

Let $j\ge2$. In $A_j$ write $r=r_1$, $b=s_1$ and $s_i=b+s'_i$ for $i\ge2$ (so $0=s'_1\le s'_2\le\dots\le s'_j$
and $s_i-s_{i-1}=s'_i-s'_{i-1}$ for $i\ge2$), and $N'_i=r_i+2s'_i$. By \eqref{eq:01} the exponent of
level $i\ge2$ is $N_i(N_i+1)/2+r_i(r_i-1)/2$ with $N_i=N'_i+2b$, and
$N_i(N_i+1)/2=N'_i(N'_i+1)/2+2bN'_i+(2b^2+b)$. Since
$(q)_R/((q)_{R-r_1}(q)_{r_1-r_2})=\qbin R{r_1}(q)_{r_1}/(q)_{r_1-r_2}$, regrouping gives
\begin{equation}\label{eq:31}
A_j(R;x)=\sum_{r=0}^{R}\sum_{b\ge0}\qbin Rr\frac{x^{r+2b}q^{(r+2b)(r+2b+1)/2+r(r-1)/2}}{(q^2;q^2)_b}\,
x^{2b(j-1)}q^{(j-1)(2b^2+b)}\,A_{j-1}(r;xq^{2b}).
\end{equation}
By the induction hypothesis $A_{j-1}(r;xq^{2b})=D_{j-1}(r;xq^{2b})$. In $D_{j-1}(r;xq^{2b})$ put
$M'_i=M_i+2b$ ($1\le i\le j-1$). Then
\begin{gather*}
x^{2b(j-1)}q^{(j-1)(2b^2+b)}\,x^{\sum M_i}q^{2b\sum M_i+\sum M_i(M_i+1)/2}=x^{\sum M'_i}q^{\sum M'_i(M'_i+1)/2},\\
(xq^{2b}q^{r+M_{j-1}+1};q)_\infty=(xq^{r+M'_{j-1}+1};q)_\infty,
\end{gather*}
the differences $M_i-M_{i+1}=M'_i-M'_{i+1}$ ($i\le j-2$) are unchanged and
$(q)_{M_{j-1}}=(q)_{M'_{j-1}-2b}$. Summing over $b$ inside (the family is summable),
\begin{multline*}
A_j(R;x)=\sum_{M'_1\ge\dots\ge M'_{j-1}\ge0}\frac{x^{\sum M'_i}q^{\sum M'_i(M'_i+1)/2}}{\prod_{i=1}^{j-2}(q)_{M'_i-M'_{i+1}}}\\
\times\sum_{r=0}^{R}\sum_{b=0}^{\lfloor M'_{j-1}/2\rfloor}\qbin Rr
\frac{x^{r+2b}q^{(r+2b)(r+2b+1)/2+r(r-1)/2}(xq^{r+M'_{j-1}+1};q)_\infty}{(q^2;q^2)_b(q)_{M'_{j-1}-2b}}.
\end{multline*}
The second line is the left side of \eqref{eq:star2} with $M=M'_{j-1}$; replacing it by the right side
and calling the summation index $L=:M_j$ (so $M'_{j-1}\ge M_j\ge0$, with denominators
$(q)_{M'_{j-1}-M_j}(q)_{M_j}$) gives exactly $D_j(R;x)$.
\end{proof}

\section{The \texorpdfstring{$x$}{x}-refinement and the proof of Theorem~\ref{thm:main}}\label{sec:x}

Let $S_k(x)$ and $\Mel_k(x)$ denote the left and the right side of \eqref{eq:xref}.

\begin{proof}[Proof of the first equality of Theorem~\ref{thm:x}]
For $k=0$ both sides are $1$. Let $k\ge1$ and $R\ge0$. The terms of $S_k(x)$ and of $A_k(R;x)$ are
indexed by the same $(r,s)$ (those with $r_1\le R$ in $A_k$); a term has $q$-order $\ge r_1^2$; and
$(q)_R/(q)_{R-r_1}=(q^{R-r_1+1};q)_{r_1}\equiv1\bmod q^{R-r_1+1}$. As $r_1^2+R-r_1+1\ge R+1$, and
$r_1^2\ge R+1$ for $r_1>R$, we get $A_k(R;x)\equiv S_k(x)\bmod q^{R+1}$. Likewise
$(xq^{R+M_k+1};q)_\infty\equiv1\bmod q^{R+1}$, so $D_k(R;x)\equiv\Mel_k(x)\bmod q^{R+1}$. By
Theorem~\ref{thm:R}, $S_k(x)\equiv\Mel_k(x)\bmod q^{R+1}$ for every $R$.
\end{proof}

\subsection{The $k$-fold sum (Bailey's lemma)}
A pair of sequences $(\alpha_n,\beta_n)_{n\ge0}$ is a \emph{Bailey pair relative to $a$} if
$\beta_n=\sum_{r=0}^n\alpha_r/((q)_{n-r}(aq;q)_{n+r})$ for all $n\ge0$. We use two classical
facts (\cite[§1 and §3]{BIS}; originally Bailey~\cite{Bailey} and Andrews~\cite{Andrews1984}; see
also~\cite{Warnaar50}):
\begin{itemize}[leftmargin=2em]
\item[(U)] $\beta_n=\delta_{n,0}$, $\alpha_n=(-1)^nq^{n(n-1)/2}(a;q)_n(1-aq^{2n})/((q)_n(1-a))$ is a Bailey
pair relative to $a$ (the unit Bailey pair);
\item[(S)] if $(\alpha,\beta)$ is a Bailey pair relative to $a$, then so is
\[\alpha'_r=a^{r/2}q^{r^2/2}\alpha_r,\qquad \beta'_n=\sum_{m=0}^{n}\frac{(-\sqrt{aq};q)_m\, a^{m/2}q^{m^2/2}\beta_m}{(q)_{n-m}(-\sqrt{aq};q)_n}\]
(the limit $\rho_1\to\infty$, $\rho_2=-\sqrt{aq}$ of Bailey's lemma).
\end{itemize}
Both are identities of rational functions in $q^{1/2}$ and $a^{1/2}$. We use them with $a=x^2q$, so
that $\sqrt{aq}=xq$, $a^{r/2}q^{r^2/2}=x^rq^{r(r+1)/2}$, $1/(aq;q)_m=1/(x^2q^2;q)_m\in\R$ and
$(a;q)_n/(1-a)=(x^2q^2;q)_{n-1}$ for $n\ge1$; so everything lies in~$\R$.

\begin{proposition}\label{prop:bailey}
For every $k\ge0$, $\Mel_k(x)$ equals \eqref{eq:melclosed} in $\R$.
\end{proposition}
\begin{proof}
Let $(\alpha^{(0)},\beta^{(0)})$ be (U) with $a=x^2q$, and let $(\alpha^{(i)},\beta^{(i)})$ be obtained by
$i$ applications of (S). Then $\alpha^{(i)}_r=x^{ir}q^{ir(r+1)/2}\alpha^{(0)}_r$, and by induction on $i\ge1$
\[
\beta^{(i)}_n=\frac1{(-xq;q)_n}\sum_{n\ge M_2\ge\dots\ge M_i\ge0}
\frac{\prod_{l=2}^{i}x^{M_l}q^{M_l(M_l+1)/2}}{(q)_{n-M_2}(q)_{M_2-M_3}\cdots(q)_{M_{i-1}-M_i}(q)_{M_i}}
\]
(for $i=1$, $\beta^{(1)}_n=1/((q)_n(-xq)_n)$; in the inductive step the factor $(-xq)_m$ of (S) cancels the
prefactor $1/(-xq)_m$ of $\beta^{(i)}_m$). The Bailey pair relation for $i=k+1$ reads
\[
\frac1{(-xq)_n}\sum_{n\ge M_1\ge\dots\ge M_k\ge0}\frac{\prod_{l=1}^kx^{M_l}q^{M_l(M_l+1)/2}}{(q)_{n-M_1}(q)_{M_1-M_2}\cdots(q)_{M_k}}
=\sum_{r=0}^n\frac{x^{(k+1)r}q^{(k+1)r(r+1)/2}\alpha^{(0)}_r}{(q)_{n-r}(x^2q^2;q)_{n+r}} .
\]
Let $n\to\infty$ in $\R$: on the left $1/((-xq)_n(q)_{n-M_1})\equiv1/((-xq)_\infty(q)_\infty)\bmod
q^{n-M_1+1}$ and the term has $q$-order $\ge M_1$; on the right
$1/((q)_{n-r}(x^2q^2)_{n+r})\equiv1/((q)_\infty(x^2q^2)_\infty)\bmod q^{n-r+1}$ and the term has
$q$-order $\ge r$; the omitted terms ($M_1>n$, $r>n$) have $q$-order $>n$. Hence
$\Mel_k(x)/((-xq)_\infty(q)_\infty)=(q)_\infty^{-1}(x^2q^2;q)_\infty^{-1}\sum_{r\ge0}x^{(k+1)r}q^{(k+1)r(r+1)/2}\alpha^{(0)}_r$,
which is \eqref{eq:melclosed} because $(x^2q^2;q)_\infty(1-x^2q)=(x^2q;q)_\infty$.
\end{proof}

This completes the proof of Theorem~\ref{thm:x}.

\begin{proof}[Proof of Theorem~\ref{thm:main}]
Put $x=1$ in Theorem~\ref{thm:x}. The left side of \eqref{eq:xref} becomes the left side of~\eqref{eq:i}.
In \eqref{eq:melclosed} the factor $(x^2q)_n/(q)_n$ becomes $1$, and
\[
\Mel_k(1)=\frac{(-q;q)_\infty}{(q;q)_\infty}\sum_{n\ge0}(-1)^nq^{((k+2)n^2+kn)/2}(1-q^{2n+1})
=\frac{(-q;q)_\infty}{(q;q)_\infty}\sum_{n\in\Z}(-1)^nq^{((k+2)n^2+kn)/2},
\]
since the terms with $-q^{2n+1}$ are the terms $n\mapsto-n-1$ of the bilateral sum
($((k+2)(n+1)^2-k(n+1))/2=((k+2)n^2+kn)/2+2n+1$). By the Jacobi triple product
$\sum_{n\in\Z}(-1)^nz^np^{n(n-1)/2}=(z,p/z,p;p)_\infty$ with $p=q^{k+2}$, $z=q^{k+1}$, the sum is
$(q,q^{k+1},q^{k+2};q^{k+2})_\infty$, and $(-q;q)_\infty=1/(q;q^2)_\infty$ by Euler.
\end{proof}

\begin{remark}\label{rem:k2}
For $k=2$, Theorem~\ref{thm:x} has a short proof: \eqref{eq:starinf}, one application of
\eqref{eq:star2}, and $R\to\infty$. An earlier version of this work proved the case $k=2$ by a chain of
fourteen Euler and Durfee moves found by computer search (the method of the first version of
this paper for $(F_2)$, $(F_3)$).
\end{remark}

\section{Corollaries}\label{sec:cor}

\begin{proof}[Proof of Corollary~\ref{cor:z1}]
For $k=0$ both sides of \eqref{eq:conj22} are $1$. For $k\ge1$, Griffin, Ono and
Warnaar~\cite[Theorem~1.2, eq.~(1.9)]{GOW} prove, for positive integers $m,n$ and $\kappa=2m+2n+2$,
\begin{multline*}
\sum_{\lambda_1\le m}q^{|\lambda|}P_{2\lambda}(1,q,q^2,\dots;q^{2n})\\=
\frac{(q^2;q^2)_\infty(q^{\kappa/2};q^{\kappa/2})_\infty(q^\kappa;q^\kappa)_\infty^{n-1}}{(q)_\infty^{n+1}}
\prod_{i=1}^n\theta(q^i;q^{\kappa/2})\prod_{1\le i<j\le n}\theta(q^{j-i},q^{i+j};q^\kappa),
\end{multline*}
For $n=1$, $m=k$ the right side is
\[\frac{(q^2;q^2)_\infty(q^{k+2};q^{k+2})_\infty\theta(q;q^{k+2})}{(q)_\infty^2}=\frac{(q,q^{k+1},q^{k+2};q^{k+2})_\infty}{(q;q^2)_\infty(q)_\infty},\]
the product side of~\eqref{eq:i}. The right side of \eqref{eq:conj22} at $z=1$ is the left side
of~\eqref{eq:i}. Theorem~\ref{thm:main} finishes the proof.
\end{proof}

The Hall--Littlewood side of \eqref{eq:conj22} at $z=1$ was also computed directly from Macdonald's
tableau formula~\cite[III.(5.11$'$)]{Macdonald} by the referees ($k\le4$, modulo $q^{26}$); it agrees with both sides. In terms of KRTW's generating functions
$\mathcal D^{(2)}_{\Lambda}(z,q)$ and $\mathcal C^{(1)}_{\Lambda}(z,q)$ of coloured partitions of types
$D^{(2)}_3$ and $C^{(1)}_1$, Theorem~\ref{thm:main} also says that the first multisum of their
Conjecture~4.9 at $z=1/q$ equals $\mathcal D^{(2)}_{k\Lambda_0}(1/q,q)=\mathcal C^{(1)}_{k\Lambda_0}(1,q)$ (their (4.7) and the
Meurman--Primc theorem), i.e.\ Conjecture~4.9 (first line) holds at $z=1/q$. For general $z$,
Conjecture~2.2 remains open: KRTW's grading $z^{\sum(r_i+s_i)}$ is not the grading
$x^{\sum(r_i+2s_i)}$ of Theorem~\ref{thm:x}.

\begin{corollary}[a second parameter]\label{cor:y}
For $k\ge1$, in $\Z[x,y][[q]]$,
\begin{multline*}
\sum_{r,s}y^{r_1}x^{\sum_i(r_i+2s_i)}\prod_{i=1}^k\frac{q^{(r_i+s_i)^2+s_i^2+s_i}}{(q)_{r_i-r_{i+1}}(q^2;q^2)_{s_i-s_{i-1}}}\\
=\sum_{M_1\ge\dots\ge M_k\ge0}\ \sum_{l\ge0}\frac{(-1)^l(y;q)_l\,x^{l+\sum_iM_i}q^{l(l+1)/2+lM_k+\sum_iM_i(M_i+1)/2}}{(q)_l\prod_{i=1}^k(q)_{M_i-M_{i+1}}} .
\end{multline*}
\end{corollary}
\begin{proof}
Work in $\Z[x][[q]][[y]]$. Multiply Theorem~\ref{thm:R} ($j=k$) by $y^R/(q)_R$ and sum over $R\ge0$. On
the left, $\sum_{R\ge r_1}y^R/(q)_{R-r_1}=y^{r_1}/(y;q)_\infty$. On the right expand
$(xq^{R+M_k+1};q)_\infty=\sum_l(-x)^lq^{l(l+1)/2+lM_k}q^{Rl}/(q)_l$ (Lemma~\ref{lem:qbin}) and use
$\sum_Ry^Rq^{Rl}/(q)_R=1/(yq^l;q)_\infty$. Multiply by $(y;q)_\infty$ and note
$(y)_\infty/(yq^l)_\infty=(y)_l$. Both sides have coefficients of $q^n$ that are polynomials in $x,y$.
\end{proof}
At $y=1$ this is Theorem~\ref{thm:x}; at $y=0$ (all $r_i=0$) it is Theorem~\ref{thm:R} with $R=0$.

\begin{remark}[a numerical observation]\label{rem:X1}
Theorem~\ref{thm:x} may be specialised at $x=1/q$ (the $q$-order minus the $x$-degree of a term is
$\ge\sum_iN_i(N_i-1)/2$, so both sides remain summable). The resulting series $S_k(1/q)$ is the value at
$z=1/q$ of the second multisum of KRTW's Conjecture~4.9. We observed, for $k\le6$ modulo $q^{80}$, that
\[
S_k(1/q)=\Mel_k(1/q)=\frac{(-q;q)_\infty}{(q;q)_\infty}\sum_{i=1}^{k+1}(q^i,q^{k+2-i},q^{k+2};q^{k+2})_\infty,
\]
the sum of the $k+1$ principally specialised level-$k$ characters of $A_1^{(1)}$; the first equality is
Theorem~\ref{thm:x}, the second is only numerical.
\end{remark}

\section{Identities (ii)--(iv) for every \texorpdfstring{$k$}{k}}\label{sec:k3}

We use $Q=q^2$ and the series $K_k$, $B_m$, $C_k$, $J_k$, $K^\Omega_k$, $B^{(k)}_{k+1}$ of §\ref{sec:intro}. Each of
them lies in $\Z[x][[q]]$: every coefficient of $q^N$ is a polynomial in $x$, because the $q$-exponents grow with
the summation indices (for $K^\Omega_k$ expand $\Omega$ into its $2k-1$ monomials, all with exponents $\ge0$). The
substitutions $x=1$, $x=q$ are therefore defined and are ring homomorphisms to $\Z[[q]]$.
§\ref{ssec:closed} recalls closed forms of $B_{k+1}$ and $J_k$ and reduces \eqref{eq:ii}--\eqref{eq:iv} to
identities in~$x$; §§\ref{ssec:transfer}--\ref{ssec:proofF} prove $(F_k)$; §\ref{ssec:iv} proves
Theorem~\ref{thm:iv}; §\ref{ssec:HL} gives the Hall--Littlewood form.

\subsection{Closed forms for all \texorpdfstring{$k$}{k}}\label{ssec:closed}
For $e\in\{1,2\}$ and $m\ge1$ put
\[
\mathfrak J^{(e)}_m(x):=\sum_{\sigma_1\ge\dots\ge\sigma_m\ge0}\frac{x^{2\sum\sigma_i}q^{2\sum\sigma_i^2}}{\prod_{i<m}(Q;Q)_{\sigma_i-\sigma_{i+1}}\cdot(Q;Q)_{\sigma_m}(-xq^e;Q)_{\sigma_m}} .
\]

\begin{theorem}\label{thm:B}
For every $k\ge1$, in $\Z[x][[q]]$:
\begin{align*}
&\text{(B1)}\quad (xq;q)_\infty B_{k+1}(x)=\sum_{n\ge0}(-1)^nx^{(k+1)n}q^{(k+2)n^2}\frac{(x^2;Q)_n(1-xq^{2n})}{(Q;Q)_n(1-x)};\\
&\text{(B2)}\quad (x^2Q;Q)_\infty J_k(x)=\sum_{m\ge0}(-1)^mx^{(k+1)m}q^{(k+2)m^2-m}\frac{(x^2;Q)_m(1-x^2q^{4m})}{(Q;Q)_m(1-x^2)};\\
&\text{(B3)}\quad B_{k+1}(x)=(-xQ;Q)_\infty\sum_{t_1\ge\dots\ge t_k\ge0}\frac{x^{\sum t_i}q^{\sum t_i^2}(-xq;Q)_{t_k}}{\prod_{i<k}(Q;Q)_{t_i-t_{i+1}}\cdot(Q;Q)_{t_k}(-xQ;Q)_{t_k}};\\
&\text{(B4)}\quad (xq;q)_\infty(x^2Q;Q)_\infty B_{k+1}(x)J_k(x)=L_k(x),\quad\text{where}\\
&\qquad L_k(x):=\tfrac12\sum_{n,m\ge0}(-1)^{n+m}x^{(k+1)(n+m)}q^{(k+2)(n^2+m^2)-n-m}\\
&\qquad\qquad\qquad\times\frac{(x^2;Q)_n(x^2;Q)_m(q^n+q^m)(1+xq^{n+m})(1-xq^{2n})(1-xq^{2m})}{(Q;Q)_n(Q;Q)_m(1-x)(1-x^2)};\\
&\text{(B5)}\quad\text{for }k=2h\text{ even, } B_{2h+1}(x)=(-xq;q)_\infty\,\mathfrak J^{(2)}_h(x)\ \text{ and }\ J_{2h}(x)=\mathfrak J^{(1)}_h(x).
\end{align*}
\end{theorem}
In (B1), (B2), (B4) the quotients $(x^2;Q)_n/(1-x)$ and $(x^2;Q)_m/(1-x^2)$ are polynomials for $n,m\ge1$,
and the terms with $n=0$ or $m=0$ are read as limits (for instance $(1-xq^{0})/(1-x)=1$), so every term
lies in $\Z[x][[q]]$.

\begin{proof}
We use Bailey pairs relative to $a$ in a base $p$: $\beta_n=\sum_{r\le n}\alpha_r/((p;p)_{n-r}(ap;p)_{n+r})$,
and the following classical facts (\cite[§1, §3]{BIS}, \cite{Andrews1984}, \cite{Warnaar50}), all
identities of rational functions in $p^{1/2}$, $a^{1/2}$: the unit pair (U) in base $p$; the limits of
Bailey's lemma
\begin{itemize}[leftmargin=2.2em]
\item[(L$\infty$)] $\alpha'_n=a^np^{n^2}\alpha_n$, $\beta'_n=\sum_{j\le n}a^jp^{j^2}\beta_j/(p;p)_{n-j}$
($\rho_1,\rho_2\to\infty$);
\item[(S)] $\alpha'_n=a^{n/2}p^{n^2/2}\alpha_n$,
$\beta'_n=\sum_{j\le n}\frac{(-\sqrt{ap};p)_ja^{j/2}p^{j^2/2}\beta_j}{(p;p)_{n-j}(-\sqrt{ap};p)_n}$
($\rho_1\to\infty$, $\rho_2=-\sqrt{ap}$);
\item[(S$'$)] $\alpha'_n=a^{n/2}p^{n(n+1)/2}\frac{1+\sqrt a}{1+\sqrt ap^n}\alpha_n$,
$\beta'_n=\sum_{j\le n}\frac{(-\sqrt a;p)_ja^{j/2}p^{j(j+1)/2}\beta_j}{(p;p)_{n-j}(-\sqrt ap;p)_n}$
($\rho_1\to\infty$, $\rho_2=-\sqrt a$);
\end{itemize}
and the limit $\lim_{n\to\infty}\beta_n=((p;p)_\infty(ap;p)_\infty)^{-1}\sum_r\alpha_r$, justified
coefficientwise exactly as in the proof of Proposition~\ref{prop:bailey}.

\emph{(B1).} Base $p=q$, $a=x$. The pair $\beta_n=1/(Q;Q)_n$ (a classical Bailey pair, cf.\ Slater's
list~\cite{Slater} and the Bailey lattice of Agarwal, Andrews and Bressoud~\cite{AAB}; (B1) is an
$x$-form of Bressoud's identities~\cite{Bressoud79,Bressoud80}),
$\alpha_n=(-1)^nq^{n^2}(a^2;Q)_n(1-aq^{2n})/((Q;Q)_n(1-a))$ is a Bailey pair relative to $a$, for any~$a$. Indeed, by Bailey
inversion
\[\alpha_n=\frac{1-aq^{2n}}{1-a}\sum_{j\le n}\frac{(a;q)_{n+j}(-1)^{n-j}q^{(n-j)(n-j-1)/2}\beta_j}{(q)_{n-j}};\]
with $\beta_j=1/((q)_j(-q)_j)$ the sum is
$\frac{(a)_n(-1)^nq^{n(n-1)/2}}{(q)_n}\,{}_2\phi_1(q^{-n},aq^n;-q;q,q)$, and the $q$-Chu--Vandermonde sum
${}_2\phi_1(q^{-n},b;c;q,q)=(c/b)_nb^n/(c)_n$ \cite[(1.5.3)]{GR} together with
$(-q^{1-n}/a)_n=(-a)_na^{-n}q^{-n(n-1)/2}$ gives the claim. Take $a=x$ and apply (L$\infty$) $k+1$ times and the limit:
the $\beta$-side is $B_{k+1}(x)/(q)_\infty$ and the $\alpha$-side is
$((q)_\infty(xq)_\infty)^{-1}\sum_nx^{(k+1)n}q^{(k+1)n^2}\alpha_n$.

\emph{(B2).} Base $Q$, $a=x^2$ (the change-of-base chain of~\cite{BIS}). Applying (S) $k+1$ times to (U), the factors $(-xq;Q)_j$ cancel step by
step; the $\beta$-limit is $C_k(x)/((Q;Q)_\infty(-xq;Q)_\infty)$ (the summation indices
$M_l=s_{k+1-l}$), the $\alpha$-side gives the right side of (B2) divided by $(Q;Q)_\infty(x^2Q;Q)_\infty$,
and $C_1(x)=(-xq;Q)_\infty$ by Euler.

\emph{(B3).} Base $Q$, $a=x^2$. (S$'$) applied to (U) gives $\beta_n=1/((Q;Q)_n(-xQ;Q)_n)$ and
$\alpha_n=(-1)^nx^nq^{2n^2}(x^2;Q)_n(1-xq^{2n})/((Q;Q)_n(1-x))$ (using
$(1+x)(1-x^2q^{4n})/((1+xq^{2n})(1-x^2))=(1-xq^{2n})/(1-x)$). Apply (S) $k$ times ($\sqrt{aQ}=xq$,
$a^{n/2}Q^{n^2/2}=x^nq^{n^2}$): $\alpha^{(k)}_n=x^{kn}q^{kn^2}\alpha_n$ is, by (B1), the $n$-th term of
$(xq;q)_\infty B_{k+1}(x)$, and
\[\beta^{(k)}_N=\frac1{(-xq;Q)_N}\sum_{N\ge t_1\ge\dots\ge t_k}\frac{x^{\sum t_i}q^{\sum t_i^2}(-xq;Q)_{t_k}}{(Q;Q)_{N-t_1}\cdots(Q;Q)_{t_{k-1}-t_k}(Q;Q)_{t_k}(-xQ;Q)_{t_k}}.\]
The limit gives (B3), because $(x^2Q;Q)_\infty/((xq;q)_\infty(-xq;Q)_\infty)=(-xQ;Q)_\infty$.

\emph{(B5).} $J_{2h}=\mathfrak J^{(1)}_h$: (S) then (L$\infty$)$^h$ in base $Q$, $a=x^2$, applied to (U),
has $\alpha_n=x^{(2h+1)n}q^{(2h+1)n^2}\alpha^{\rm U}_n$, the $n$-th term of (B2), and $\beta$-limit
$\mathfrak J^{(1)}_h/(Q;Q)_\infty$. $B_{2h+1}=(-xq;q)_\infty\mathfrak J^{(2)}_h$: (S$'$) then
(L$\infty$)$^h$ has $\alpha_n=x^{2hn}q^{2hn^2}$ times the $\alpha$ of (B3), the $n$-th term of (B1) with
$k=2h$, and $\beta$-limit $\mathfrak J^{(2)}_h/(Q;Q)_\infty$; finally
$(x^2Q;Q)_\infty/(xq;q)_\infty=(-xq;q)_\infty$.

\emph{(B4).} Let $b_n$, $j_m$ be the $n$-th and $m$-th terms of (B1), (B2). Then the left side of (B4) is
$\sum_{n,m}b_nj_m=\frac12\sum_{n,m}(b_nj_m+b_mj_n)$, and
\[b_nj_m+b_mj_n=(-1)^{n+m}x^{(k+1)(n+m)}q^{(k+2)(n^2+m^2)}\frac{(x^2;Q)_n(x^2;Q)_m}{(Q;Q)_n(Q;Q)_m(1-x)(1-x^2)}\,E\]
with $E=q^{-m}(1-xq^{2n})(1-x^2q^{4m})+q^{-n}(1-xq^{2m})(1-x^2q^{4n})$. With $u=q^n$, $v=q^m$ one checks by
expanding that $uvE=(u+v)(1+xuv)(1-xu^2)(1-xv^2)$, which gives $L_k(x)$.
\end{proof}

Put $\varpi_n:=x^nq^{n^2}$ and
\begin{equation}\label{eq:betas}
\begin{gathered}
\beta^B_n:=\frac{(-1)^nq^{n^2}(x^2;Q)_n(1-xq^{2n})}{(Q;Q)_n(1-x)},\qquad
\beta^J_m:=\frac{(-1)^mq^{m^2-m}(x^2;Q)_m(1-x^2q^{4m})}{(Q;Q)_m(1-x^2)},\\
b^\Omega_n:=\beta^B_n\,\frac{q^{-2n}+x-1+x^2q^{2n}}{x(1+x)},
\end{gathered}
\end{equation}
so that (B1), (B2) read $(xq;q)_\infty B_{k+1}=\sum_n\varpi_n^{k+1}\beta^B_n$ and
$(x^2Q;Q)_\infty J_k=\sum_m\varpi_m^{k+1}\beta^J_m$. One has $b^\Omega_0=1$ and, for $n\ge1$ (use
$(x^2;Q)_n/(1-x)=(1+x)(x^2Q;Q)_{n-1}$),
\begin{equation}\label{eq:bOmega}
b^\Omega_n=\frac{(-1)^nq^{n^2-2n}(x^2Q;Q)_{n-1}(1-xq^{2n})\bigl(1-(1-x)q^{2n}+x^2q^{4n}\bigr)}{x\,(Q;Q)_n},
\end{equation}
so $\varpi_n^{k+1}b^\Omega_n$ is $x^{(k+1)n-1}q^{(k+2)n^2-2n}$ times a polynomial in $x$ with coefficients in $\Z[[q]]$.

\begin{proposition}\label{prop:Fimplies}
\begin{enumerate}[label=(\alph*)]
\item If $(F_k)$ holds for some $k\ge1$, then \eqref{eq:ii} and \eqref{eq:iii} hold for that $k$.
\item If $(xq;q)_\infty K^\Omega_k(x)=J_k(x)\sum_{n\ge0}\varpi_n^{k+1}b^\Omega_n$ holds in $\Z[x][[q]]$ for some $k\ge1$,
then \eqref{eq:iv} holds for that $k$.
\end{enumerate}
\end{proposition}
\begin{proof}
(a) At $x=1$ and $x=q$ the sums in (B1), (B2) are bilateral theta series; we write
$\vartheta(z;p):=(p;p)_\infty\theta(z;p)=(z,p/z,p;p)_\infty$. At $x=1$ the factor $(x^2;Q)_n(1-xq^{2n})/((Q;Q)_n(1-x))$ becomes $2$ for
$n\ge1$ and $1$ for $n=0$, so $(q)_\infty B_{k+1}(1)=\sum_{n\in\Z}(-1)^nq^{(k+2)n^2}=\vartheta(q^{k+2};q^{2k+4})$; at
$x=q$ it becomes $(1-q^{2n+1})/(1-q)$, $(xq;q)_\infty=(q)_\infty/(1-q)$, and the terms $n$ and $-n-1$ combine
to $(q)_\infty B_{k+1}(q)=\sum_{n\in\Z}(-1)^nq^{(k+2)n^2+(k+1)n}=\vartheta(q;q^{2k+4})$. In (B2), at $x=1$ the
factor $(x^2;Q)_m(1-x^2q^{4m})/((Q;Q)_m(1-x^2))$ is $1+q^{2m}$ for $m\ge1$ (and $1$ for $m=0$), giving
$(Q;Q)_\infty J_k(1)=\sum_{m\in\Z}(-1)^mq^{(k+2)m^2-m}=\vartheta(q^{k+1};q^{2k+4})$; at $x=q$ it is
$(1-q^{4m+2})/(1-q^2)$ and $m$, $-m-1$ combine to $(Q;Q)_\infty J_k(q)=\vartheta(q^2;q^{2k+4})$ (Jacobi
triple product throughout). Hence, under $(F_k)$,
$K_k(1)=\vartheta(q^{k+2};q^{2k+4})\vartheta(q^{k+1};q^{2k+4})/((q)_\infty(Q;Q)_\infty)$ and
$K_k(q)=\vartheta(q;q^{2k+4})\vartheta(q^2;q^{2k+4})/((q)_\infty(Q;Q)_\infty)$, which are the right sides of
\eqref{eq:iii} and~\eqref{eq:ii}.
(b) $K^\Omega_k(1)$ is the left side of~\eqref{eq:iv}. By~\eqref{eq:bOmega}, at $x=1$ we have $b^\Omega_0=1$ and
$b^\Omega_n=(-1)^nq^{n^2-2n}(1+q^{4n})$ for $n\ge1$, so
\[
\sum_{n\ge0}\varpi_n^{k+1}b^\Omega_n\Big|_{x=1}=1+\sum_{n\ge1}(-1)^n\bigl(q^{(k+2)n^2-2n}+q^{(k+2)n^2+2n}\bigr)
=\sum_{n\in\Z}(-1)^nq^{(k+2)n^2-2n}=\vartheta(q^k;q^{2k+4})
\]
by the Jacobi triple product. With $(xq;q)_\infty|_{x=1}=(q)_\infty$ and $J_k(1)$ from (a),
$K^\Omega_k(1)=\vartheta(q^k;q^{2k+4})\vartheta(q^{k+1};q^{2k+4})/((q)_\infty(Q;Q)_\infty)$, the right side
of~\eqref{eq:iv}.
\end{proof}

\begin{remark}\label{rem:k1}
For $k=1$, $J_1=1$, and $K_1(x)=B_2(x)$ holds term by term (put $N_2=s_1$, $N_1=r_1+s_1$). So
$(F_1)$ is trivial, and \eqref{eq:ii}, \eqref{eq:iii} for $k=1$ follow from (B1); the evaluation of
$B_2(1)$ is an instance of Bressoud's identities. Likewise $\Omega=q^{r_1+2s_1}$ for $k=1$, so
$K^\Omega_1(x)=K_1(xq)=B_2(xq)=B^{(1)}_2(x)$.
\end{remark}

The rest of this section proves the two hypotheses of Proposition~\ref{prop:Fimplies} for every $k$.

\subsection{The transfer operator}\label{ssec:transfer}
\emph{Ring.} In §§\ref{ssec:transfer}--\ref{ssec:iv} we work in $\mathbb F:=\Q(q)((x))$, the field of formal
Laurent series in $x$ over $\Q(q)$, with the $x$-adic valuation $\operatorname{ord}$; rational functions of $x$ and
$q$ are elements of $\mathbb F$ (expanded in $x$). For each $D$, the coefficient of $x^D$ in each of the series of
§\ref{sec:intro} and of Theorem~\ref{thm:B} is a finite sum of rational functions of $q$ that are regular at $q=0$
(the exponent of $x$ bounds the summation indices; $C_1\equiv1\bmod x$, which is used only to define $J_k$). The
coefficientwise Taylor map from such series to $\Q[[q]][[x]]$ is an injective ring homomorphism, and
$\Z[x][[q]]\subset\Q[[x,q]]=\Q[[q]][[x]]$. Hence identities between these series may be proved in~$\mathbb F$, and
(B1), (B2), proved in $\Z[x][[q]]$, hold in~$\mathbb F$. Infinite products such as $(xq;q)_\infty$, $(x^2Q;Q)_\infty$ and
$1/(aq;q)_\infty$ ($a=xq^{2s}$) do not converge $x$-adically in~$\mathbb F$ ($\Q(q)$ carries the discrete topology); they
denote the elements of $\mathbb F$ given by their Euler expansions, e.g.\
$(xq;q)_\infty=\sum_n(-1)^nq^{n(n+1)/2}x^n/(q)_n$ and $1/(zq;q)_\infty=\sum_nz^nq^n/(q)_n$, whose $x$-coefficients are rational
functions of $q$ regular at $q=0$; every identity used between such products and the series above (for instance
(B-ii) below) holds in $\mathbb F$ by the same transfer.

\emph{Functions.} Let $\mathcal V$ be the space of functions $F\colon\N^2\to\mathbb F$ with
$\inf\operatorname{ord}F>-\infty$, extended by $F(R,S):=0$ when $R<0$ or $S<0$. Put
$w(r,s):=x^{r+2s}q^{(r+s)^2+s^2}$ (so $\operatorname{ord}w(r,s)=r+2s$) and
\[
(TF)(R,S):=\sum_{r=0}^R\sum_{s\ge S}\frac{w(r,s)F(r,s)}{(q)_{R-r}(Q;Q)_{s-S}},\qquad u(R,S):=\frac1{(q)_R},\qquad
\omega(F):=\sum_{r,s\ge0}\frac{w(r,s)F(r,s)}{(Q;Q)_s}.
\]
The sums converge $x$-adically for $F\in\mathcal V$, and $T\mathcal V\subset\mathcal V$. Unwinding the definitions
(level $i$ of $K_k$ is the pair $(r_i,s_i)$; the factor $1/(q)_{r_k}$ is $u$, and the factor $1/(Q;Q)_{s_1}$ is
in~$\omega$),
\begin{equation}\label{eq:Kop}
K_k(x)=\omega(T^{k-1}u)\qquad(k\ge1).
\end{equation}

\emph{Operators.} $P$ and $Y$ are the multiplications by $q^R$ and $q^{2S}$; $\mathsf SF(R,S):=F(R-1,S)$,
$\mathsf UF(R,S):=F(R,S+1)$; $\mathcal W$ is the multiplication by $w$. Then $\mathsf SP=q^{-1}P\mathsf S$,
$\mathsf UY=QY\mathsf U$, all other pairs among $P,Y,\mathsf S,\mathsf U$ commute, and
\begin{equation}\label{eq:TEW}
T=E\mathcal W,\quad E:=\sum_{j,b\ge0}\frac{\mathsf S^j\mathsf U^b}{(q)_j(Q;Q)_b},\qquad
\mathcal W\mathsf S=xq^{-1}P^2Y\mathsf S\mathcal W,\quad \mathcal W\mathsf U=x^{-2}q^{-2}P^{-2}Y^{-2}\mathsf U\mathcal W
\end{equation}
(because $w(R,S)/w(R-1,S)=xq^{2R+2S-1}$ and $w(R,S)/w(R,S+1)=x^{-2}q^{-2R-4S-2}$).

\subsection{Two commuting difference operators}\label{ssec:comm}
Let
\begin{align*}
H'&:=(Y^{-1}+x)P^{-1}(1-\mathsf S)+xP(1+xY)-x^2PY\mathsf U,\\
G&:=P^{-2}Y^{-1}(1-\mathsf S)(1-q\mathsf S)+Y^{-1}-x(1+q^{-1})\mathsf S+x^2(1+P^2)Y(1-\mathsf U).
\end{align*}

\begin{proposition}\label{prop:comm}
$TH'=H'T$ and $TG=GT$ on $\mathcal V$.
\end{proposition}
\begin{proof}
Every operator below is a series $\sum_{j,b}c_{j,b}(P,Y)\mathsf S^j\mathsf U^b$ with Laurent-polynomial coefficients
written to the left (``normal order''). They are applied to the class of functions $F$ on $\Z^2$ supported in
$\{R\ge0,\ S\ge-1\}$ with $\operatorname{ord}F(R,S)\ge R+2S+c$ for some constant $c$ (it contains $\mathcal WF$ and $H'\mathcal WF$
for $F\in\mathcal V$); on it such a series converges, because $F(R,S+b)$ has order $\ge2b+R+2S+c$, and all
rearrangements below are finite recombinations of $x$-adically convergent sums. Products are normal-ordered with
$\mathsf S^jP^m=q^{-jm}P^m\mathsf S^j$, $\mathsf U^bY^n=Q^{bn}Y^n\mathsf U^b$, and equal normal-ordered coefficients give
equal operators on $\mathcal W\mathcal V$. We compute with functions on $\Z^2$ that vanish outside $\N^2$:
$\mathsf U$, $E$ and $H'$ can produce values off $\N^2$ (e.g.\ $(H'F)(R,-1)=-x^2q^{R-2}F(R,0)$), but the restriction
of $E\mathcal W$ to $\N^2$ is $T$, it reads only $R'\le R$ and $S'\ge S$, and $E\mathcal WF$ vanishes for $R<0$; so
an identity $E\mathcal WH'=H'E\mathcal W$ on $\Z^2$, restricted to $\N^2$, is $TH'=H'T$ on $\mathcal V$ (likewise for
$G$ and for the operators of §\ref{ssec:iv}).
By~\eqref{eq:TEW}, $\mathcal WH'=\tilde H'\mathcal W$ with
\[
\tilde H'=(Y^{-1}+x)P^{-1}-xq^{-1}P(1+xY)\mathsf S+xP(1+xY)-q^{-2}P^{-1}Y^{-1}\mathsf U .
\]
It suffices to show $H'E=E\tilde H'$; then $TH'=E\mathcal WH'=E\tilde H'\mathcal W=H'E\mathcal W=H'T$. Using
$1/(q)_j-1/(q)_{j-1}=q^j/(q)_j$ and $1/(Q;Q)_{b-1}=(1-Q^b)/(Q;Q)_b$, the coefficient of $\mathsf S^j\mathsf U^b$,
multiplied by $(q)_j(Q;Q)_b$, is
\begin{align*}
\text{in }H'E:&\quad (Y^{-1}+x)q^jP^{-1}+xP(1+xY)-x^2PY(1-Q^b),\\
\text{in }E\tilde H':&\quad (Q^{-b}Y^{-1}+x)q^jP^{-1}-xq^{-j}(1-q^j)P(1+xQ^bY)\\
&\qquad+xq^{-j}P(1+xQ^bY)-q^jQ^{-b}(1-Q^b)P^{-1}Y^{-1},
\end{align*}
and both equal $q^jP^{-1}Y^{-1}+xq^jP^{-1}+xP+x^2Q^bPY$. For $G$: $\mathcal WG=\tilde G\mathcal W$ with
\begin{multline*}
\tilde G=P^{-2}Y^{-1}+Y^{-1}-x(1+q^{-1})\mathsf S-x^2q^{-1}(1+q^{-1})P^2Y\mathsf S\\
+x^2q^{-3}P^2Y\mathsf S^2+x^2(1+P^2)Y-q^{-2}(1+P^{-2})Y^{-1}\mathsf U,
\end{multline*}
and the coefficients of $\mathsf S^j\mathsf U^b$ (times $(q)_j(Q;Q)_b$) of $GE$ and $E\tilde G$ both equal
$q^{2j}P^{-2}Y^{-1}+Y^{-1}-x(1+q^{-1})(1-q^j)+x^2Q^b(1+P^2)Y$ (for $GE$ use
$1-(1+q)(1-q^j)+q(1-q^j)(1-q^{j-1})=q^{2j}$).
\end{proof}

With the rank-one operators $h_a:=P^{-1}(1-\mathsf S)+aP$ (acting in $R$; $a$ may depend on $S$) one has
\begin{equation}\label{eq:HGh}
H'=(Y^{-1}+x)h_{xY}-x^2PY\mathsf U,\qquad G=Y^{-1}h_{xY}^2+Y^{-1}(1-xY)^2-x^2(1+P^2)Y\mathsf U.
\end{equation}

\subsection{The Bailey basis and the eigenvectors of \texorpdfstring{$T$}{T}}\label{ssec:basis}
For $a$ with all $(aq)_n$ invertible put $\varphi_e(R;a):=1/((q)_{R-e}(aq)_{R+e})$ ($=0$ for $R<e$). We use:
\begin{enumerate}[label=(B-\roman*),leftmargin=3em]
\item $\sum_{r\le R}a^rq^{r^2}\varphi_e(r;a)/(q)_{R-r}=a^eq^{e^2}\varphi_e(R;a)$ (Bailey's lemma, case (L$\infty$)).
\item $\sum_{r\ge0}a^rq^{r^2}\varphi_e(r;a)=a^eq^{e^2}/(aq;q)_\infty$ (Cauchy: $\sum_jz^jq^{j^2}/((q)_j(zq)_j)=1/(zq)_\infty$,
$z=aq^{2e}$).
\item $h_a\varphi_e(\cdot;a)=(q^{-e}+aq^e)\varphi_e(\cdot;a)$, from $\varphi_e(R-1;a)/\varphi_e(R;a)=(1-q^{R-e})(1-aq^{R+e})$
(also at $R=e$, where $\varphi_e(e-1;a)=0$).
\item (Change of $a$.) $\varphi_d(\cdot;aq^c)=\sum_{e=d}^{d+c}\Gamma^{(c)}_{d,e}(a)\varphi_e(\cdot;a)$ with
\[\Gamma^{(c)}_{d,e}(a)=\frac{1-aq^{2e}}{1-a}\,\frac{(a)_{d+e}}{(aq^{c+1})_{d+e}}\,(-1)^{e-d}q^{\binom{e-d}2}a^{e-d}q^{(e+d)(e-d)}\qbin{c}{e-d}\]
(Bailey inversion and $q$-Chu--Vandermonde), and $\sum_e\Gamma^{(c)}_{d,e}(a)=(aq;q)_c$. Support and sum are
algebraic: over the common denominator $(q)_{R-d}(aq)_{R+d+c}$ the numerator of $\varphi_d(R;aq^c)$ is the constant
$(aq)_c$, while the numerators $p_e(z)=(q^{R-e+1})_{e-d}(aq^{R+e+1})_{d+c-e}$ of $\varphi_e(R;a)$, $d\le e\le d+c$,
are $c+1$ polynomials of degree $\le c$ in $z=q^R$, linearly independent ($p_e(q^m)=0$ for $d\le m<e$,
$p_e(q^e)\ne0$), hence a basis; all $p_e(0)=1$, so the coefficient sum is the value $(aq)_c$ at $z=0$. Only the support,
the coefficient sum and the case $c=1$ are used below; for $c=1$, $1-aq=A(1-aq^{d+1}z)+B(1-q^{-d}z)$ gives
$\Gamma^{(1)}_{d,d}=A=(1-aq)/(1-aq^{2d+1})$ and $\Gamma^{(1)}_{d,d+1}=B=-aq^{2d+1}A$. Also
$\varphi_d(R;aq^2)=(aq)_2\varphi_{d+1}(R+1;a)$.
\item For $n\ge1$, with $(\tau F)(R):=q^RF(R+1)$: $\tau\varphi_n=\alpha_n\varphi_{n-1}+\beta_n\varphi_n+\gamma_n\varphi_{n+1}$, where
\[\alpha_n=\frac{q^{n-1}}{(1-aq^{2n-1})(1-aq^{2n})},\quad \beta_n=\frac{-q^{n-1}(1+aq^{2n})}{(1-aq^{2n-1})(1-aq^{2n+1})},\quad
\gamma_n=\frac{aq^{3n}}{(1-aq^{2n})(1-aq^{2n+1})},\]
and $\alpha_n+\beta_n+\gamma_n=0$; likewise $\varphi_n(R+1)$ and $q^{2R}\varphi_n(R+1)$ are combinations of
$\varphi_{n-1},\varphi_n,\varphi_{n+1}$ with coefficient sums $1$ and $0$ (partial fractions in $z=q^R$; for $n=0$
these fail, e.g.\ $\tau\varphi_0$ has infinite support, and $n\ge1$ in every use below).
\end{enumerate}
For a finite combination $F=\sum_e\kappa_e\varphi_e(\cdot;a)$ put $v_a(F):=\sum_e\kappa_e$ (the $\varphi_e(\cdot;a)$ are
linearly independent, since $\varphi_e(R;a)$ vanishes for $R<e$ and not at $R=e$; for $|q|,|a|<1$, $v_a(F)$ is
$\lim_{R\to\infty}(q)_\infty(aq)_\infty F(R)$). Then $v_a(h_aF)=\sum_e(q^{-e}+aq^e)\kappa_e$; by (B-v), if
$\kappa_0=0$ then $v_a(\tau F)=0$, $v_a(F(\cdot+1))=v_a(F)$, $v_a(q^{2R}F(\cdot+1))=0$, and since
$h_a(\tau F)=F(\cdot+1)+aq^{2R}F(\cdot+1)-q^{-1}F$ (pointwise for $R\ge0$, using $F(0)=0$),
$v_a(h_a\tau F)=(1-q^{-1})v_a(F)$; by (B-iv), $v_a(\varphi_d(\cdot;aq^c))=(aq;q)_c$.

\emph{Basis of $\mathcal V$.} Put $\Xi_{e,\sigma}(R,S):=\delta_{S,\sigma}\varphi_e(R;xq^{2\sigma})$. Every $F\in\mathcal V$ is
uniquely a locally finite sum $\sum_{e,\sigma}f_{e,\sigma}\Xi_{e,\sigma}$ (Bailey inversion in each slice). Since
$w(r,\sigma)=x^{2\sigma}q^{2\sigma^2}(xq^{2\sigma})^rq^{r^2}$, (B-i) and (B-iv) give
\begin{equation}\label{eq:TXi}
T\,\Xi_{e,\sigma}=w(e,\sigma)\sum_{S=0}^{\sigma}\frac1{(Q;Q)_{\sigma-S}}\sum_{e'=e}^{e+2(\sigma-S)}\Gamma^{(2\sigma-2S)}_{e,e'}(xq^{2S})\,\Xi_{e',S}.
\end{equation}
With $(N,M):=(e+\sigma,\sigma)$ the diagonal entry is $w(e,\sigma)=\varpi_{N,M}:=x^{N+M}q^{N^2+M^2}$; $(e,\sigma)\mapsto(N,M)$
is a bijection onto $\{N\ge M\ge0\}$, and the $\varpi_{N,M}$ are pairwise distinct monomials. Write $L:=N-M$.

\begin{proposition}\label{prop:eig}
For $N\ge M\ge0$ there is a unique $\varphi_{N,M}\in\mathcal V$ with $T\varphi_{N,M}=\varpi_{N,M}\varphi_{N,M}$ of the
form $\Xi_{L,M}+\sum_{S<M}\sum_ec^{(S)}_e\Xi_{e,S}$; in it $c^{(S)}_e\ne0$ only for $L\le e\le L+2(M-S)$. Every
eigenvector of $T$ with finitely many nonzero $\Xi$-coefficients is a multiple of some $\varphi_{N,M}$. Moreover
$\varphi_{N,M}(R,S)=0$ unless $R\ge L$ and $S\le M$.
\end{proposition}
\begin{proof}
By~\eqref{eq:TXi} the eigen-equation is triangular in $S$ (downwards): for $S<M$,
$(\varpi_{N,M}-\varpi_{e+S,S})c^{(S)}_e$ is determined by the coefficients on the slices above $S$, and
$\varpi_{e+S,S}\ne\varpi_{N,M}$. The support follows from that of $\Gamma$. For an eigenvector with finitely many
$\Xi$-coefficients look at its highest slice $\sigma$: there $(\varpi-w(e,\sigma))f_{e,\sigma}=0$, so
$(e,\sigma)=(L,M)$; the eigenvector divided by $f_{L,M}$ is then of the form above, so by uniqueness it is $\varphi_{N,M}$.
\end{proof}

Since $w(r,s+1)=x^2q^2\cdot w|_{x\to xq^2}(r,s)$, the function $(R,S)\mapsto\varphi_{N,M}(R,S+1;x)$ is an eigenvector of
$T|_{x\to xq^2}$ with eigenvalue $\varpi_{N,M}(x)/(x^2q^2)=\varpi_{N-1,M-1}(xq^2)$ and top $\Xi_{L,M-1}|_{x\to xq^2}$; by
uniqueness,
\begin{equation}\label{eq:shift}
\varphi_{N,M}(\cdot,S+1;x)=\varphi_{N-1,M-1}(\cdot,S;xq^2)\qquad(M\ge1).
\end{equation}

\begin{lemma}\label{lem:HGeig}
$H'\varphi_{N,M}=\varepsilon_{N,M}\varphi_{N,M}$ and $G\varphi_{N,M}=g_{N,M}\varphi_{N,M}$, where
\[\varepsilon_{N,M}=q^{-N-M}(1+xq^{2N})(1+xq^{2M}),\qquad g_{N,M}=q^{-2N}+x^2q^{2N}+q^{-2M}+x^2q^{2M}.\]
\end{lemma}
\begin{proof}
By~\eqref{eq:HGh}, (B-iii), (B-iv) and (B-v) (note $q^R\varphi_e(R;xq^{2S+2})=(xq^{2S+1})_2(\tau\varphi_{e+1})(R)$ relative to
$xq^{2S}$, and $e+1\ge1$), $H'$ and $G$ map functions with finitely many $\Xi$-coefficients to such functions and do
not raise the top slice. By Proposition~\ref{prop:comm}, $H'\varphi_{N,M}$ is an eigenvector of $T$ for
$\varpi_{N,M}$, hence a multiple of $\varphi_{N,M}$; the multiple is read off on the slice $S=M$, where the
$\mathsf U$-terms vanish: $(q^{-2M}+x)(q^{-L}+xq^{2M+L})=\varepsilon_{N,M}$, and for $G$,
$q^{-2M}(q^{-L}+xq^{2M+L})^2+q^{-2M}(1-xq^{2M})^2=g_{N,M}$.
\end{proof}

\subsection{Values at infinity}\label{ssec:vinf}
Put $A_{N,M}(x):=v_x(\varphi_{N,M}(\cdot,0))$.

\begin{proposition}\label{prop:Arec}
$A_{L,0}=1$ and, for $M\ge1$,
\[
A_{N,M}(x)=\frac{x^2(1+x)(x+q^{-1})\,q^{2N+2M}\,(xq;q)_2}{(1-q^{2N})(1-q^{2M})(1-x^2q^{2N})(1-x^2q^{2M})}\,A_{N-1,M-1}(xq^2).
\]
\end{proposition}
\begin{proof}
$\varphi_{L,0}=\Xi_{L,0}$. Let $M\ge1$, $F_0:=\varphi_{N,M}(\cdot,0)$ (a combination of $\varphi_e(\cdot;x)$, $e\ge L$)
and $F_1:=\varphi_{N,M}(\cdot,1)$. By~\eqref{eq:shift}, $F_1=\sum_e\kappa'_e\varphi_e(\cdot;xq^2)$ with
$\sum_e\kappa'_e=A_{N-1,M-1}(xq^2)$. By (B-iv), $F_1(R)=(xq)_2\,\Phi(R+1)$ with $\Phi:=\sum_e\kappa'_e\varphi_{e+1}(\cdot;x)$, a
combination of $\varphi_n(\cdot;x)$ with $n\ge1$ and $v_x(\Phi)=\sum_e\kappa'_e$. Put $v_1:=(xq)_2A_{N-1,M-1}(xq^2)$. By the
rules of §\ref{ssec:basis} applied to $\Phi$: $q^RF_1(R)=(xq)_2(\tau\Phi)(R)$ has $v_x=0$, $h_x(q^RF_1)$ has
$v_x=(1-q^{-1})v_1$, and $(1+q^{2R})F_1(R)=(xq)_2(1+q^{2R})\Phi(R+1)$ has $v_x=v_1$. By~\eqref{eq:HGh} ($Y=1$ on $S=0$)
the slices $S=0$ of $H'\varphi=\varepsilon\varphi$ and $G\varphi=g\varphi$ read
\[
(1+x)h_xF_0-x^2q^RF_1(R)=\varepsilon F_0,\qquad h_x^2F_0+(1-x)^2F_0-x^2(1+q^{2R})F_1(R)=gF_0 .
\]
Apply $v_x$ to both equations and to $h_x$ applied to the first; with $\mu_j:=v_x(h_x^jF_0)$ ($h_x$ is diagonal on $F_0$):
\[
(1+x)\mu_1-\varepsilon\mu_0=0,\qquad (1+x)\mu_2-\varepsilon\mu_1=x^2(1-q^{-1})v_1,\qquad \mu_2+(1-x)^2\mu_0-x^2v_1=g\mu_0 .
\]
Eliminating $\mu_1,\mu_2$ gives $\mu_0[\varepsilon^2+(1+x)^2((1-x)^2-g)]=x^2(1+x)(x+q^{-1})v_1$. With $e_n:=q^{-n}+xq^n$ one has
$\varepsilon=e_Ne_M$ and $g=e_N^2+e_M^2-4x$, so the bracket is $(e_N^2-(1+x)^2)(e_M^2-(1+x)^2)$, and
$e_n^2-(1+x)^2=q^{-2n}(1-q^{2n})(1-x^2q^{2n})$.
\end{proof}

\begin{corollary}\label{cor:V}
$A_{N,M}(x)=(xq;q)_{2M}V_{N,M}(x)$ with
\[
V_{N,M}(x):=\frac{x^{2M}q^{2MN+M(2M-1)}(-x;q)_M}{(Q;Q)_M\,(Q^{N-M+1};Q)_M\,(xq^M;q)_M\,(x^2q^{2N};Q)_M}.
\]
\end{corollary}
\begin{proof}
Iterate Proposition~\ref{prop:Arec} and use $(x^2q^{2M};Q)_M=(xq^M;q)_M(-xq^M;q)_M$.
\end{proof}

\begin{lemma}\label{lem:omphi}
$\omega(\varphi_{N,M})=\varpi_{N,M}A_{N,M}/(xq;q)_\infty$. More generally, for every finite combination $F$ of the
$\Xi_{e,\sigma}$, $(TF)(\cdot,0)$ is a finite combination of the $\varphi_e(\cdot;x)$ and $(xq;q)_\infty\,\omega(F)=v_x((TF)(\cdot,0))$.
\end{lemma}
\begin{proof}
By linearity it suffices to take $F=\Xi_{e,s}$. By (B-ii), $(xq)_\infty\omega(\Xi_{e,s})=w(e,s)(xq)_{2s}/(Q;Q)_s$; by (B-i),
$(T\Xi_{e,s})(\cdot,0)=w(e,s)\varphi_e(\cdot;xq^{2s})/(Q;Q)_s$, whose $v_x$ is the same by (B-iv). For $F=\varphi_{N,M}$ use
$(T\varphi)(\cdot,0)=\varpi_{N,M}\varphi(\cdot,0)$.
\end{proof}

\subsection{The reversed chain}\label{ssec:rev}
Let
\[
(T^\sharp\Phi)(r,s):=\sum_{R\ge r}\sum_{S\le s}\frac{w(R,S)\Phi(R,S)}{(q)_{R-r}(Q;Q)_{s-S}}
\]
($T^\sharp=\mathcal W^{-1}T^{\mathsf t}\mathcal W$ as kernels). If $T^\sharp\Phi=\varpi\Phi$, then
$\psi_\Phi(F):=\sum_{r,s}w(r,s)\Phi(r,s)F(r,s)$ satisfies $\psi_\Phi(TF)=\varpi\,\psi_\Phi(F)$ whenever the double sum
is finite (the case below). Let $\varphi^Q_e(S;a):=1/((Q;Q)_{S-e}(aQ;Q)_{S+e})$ and
$\Xi^\sharp_{e,\rho}(R,S):=\delta_{R,\rho}\varphi^Q_e(S;x^2q^{2\rho})$. Since $w(\rho,S)=x^\rho q^{\rho^2}(x^2q^{2\rho})^SQ^{S^2}$, the base-$Q$
versions of (B-i), (B-iv) give
\[
T^\sharp\Xi^\sharp_{e,\rho}=\varpi_{\rho+e,e}\sum_{r=0}^{\rho}\frac1{(q)_{\rho-r}}\sum_{e'=e}^{e+\rho-r}\Gamma^{Q,(\rho-r)}_{e,e'}(x^2q^{2r})\,\Xi^\sharp_{e',r},
\]
where $\Gamma^{Q,(c)}$ is $\Gamma^{(c)}$ with $q\to Q$ (change $a\to aQ^c$ in base $Q$; $\sum_{e'}\Gamma^{Q,(c)}_{e,e'}(a)=(aQ;Q)_c$).
Exactly as in Proposition~\ref{prop:eig} there is a unique eigenvector
$\varphi^\sharp_{N,M}=\Xi^\sharp_{M,L}+\sum_{\rho<L}\sum_ec_{e,\rho}\Xi^\sharp_{e,\rho}$ of $T^\sharp$ for $\varpi_{N,M}$, with
$c_{e,\rho}\ne0$ only for $M\le e\le M+L-\rho$; $\varphi^\sharp_{N,M}(R,S)=0$ unless $R\le L$ and $S\ge M$; and every eigenvector of
$T^\sharp$ with finitely many $\Xi^\sharp$-coefficients is a multiple of some $\varphi^\sharp_{N,M}$.

Transposing Proposition~\ref{prop:comm}, $H'^\sharp:=\mathcal W^{-1}H'^{\mathsf t}\mathcal W$ and
$G^\sharp:=\mathcal W^{-1}G^{\mathsf t}\mathcal W$ commute with $T^\sharp$ (the kernel identity
$(T^\sharp H'^\sharp)(a,b)=w(a)^{-1}w(b)(H'T)(b,a)$ involves only finite sums, and applied to $\varphi^\sharp$ all sums are
finite: $H'^\sharp$ is local, $T^\sharp$ sums over finitely many $S$, and $\varphi^\sharp$ vanishes for $R>L$). Explicitly,
with $\mathsf S^\natural F(R,S):=F(R+1,S)$, $\mathsf U^\natural F(R,S):=F(R,S-1)$ and $h^Q_a:=Y^{-1}(1-\mathsf U^\natural)+aY$ (acting
in $S$),
\begin{equation}\label{eq:HGsharp}
\begin{aligned}
H'^\sharp&=P^{-1}Y^{-1}(1-\mathsf U^\natural)+xP^{-1}+xP(1+xY)(1-\mathsf S^\natural),\\
G^\sharp&=(P^{-2}+1)h^Q_{x^2P^2}-x(1+q^{-1})\mathsf S^\natural-x^2(1+q)P^2Y\mathsf S^\natural+qx^2P^2Y(\mathsf S^\natural)^2 .
\end{aligned}
\end{equation}
The base-$Q$ facts used are: $h^Q_a\varphi^Q_e(\cdot;a)=(Q^{-e}+aQ^e)\varphi^Q_e(\cdot;a)$;
$\varphi^Q_e(S;aQ)=c_1\varphi^Q_e(S;a)+c_2\varphi^Q_{e+1}(S;a)$ with $c_1=(1-aQ)/(1-aQ^{2e+1})$,
$c_2=-(1-aQ)aQ^{2e+1}/(1-aQ^{2e+1})$, $c_1+c_2=1-aQ$; $Q^S\varphi^Q_e(S;aQ)=b_e(\varphi^Q_e(S;a)-\varphi^Q_{e+1}(S;a))$ with
$b_e=(1-aQ)/(Q^{-e}-aQ^{e+1})$; and $Q^S\varphi^Q_e(S;aQ^2)$ is a combination of $\varphi^Q_e,\varphi^Q_{e+1},\varphi^Q_{e+2}$
(relative to $a$) with coefficient sum $0$ (partial fractions in $Q^S$). With $v^Q_a$ defined as $v_a$ (coefficient
sum relative to $a$) these give, as in Lemma~\ref{lem:HGeig}, $H'^\sharp\varphi^\sharp_{N,M}=\varepsilon_{N,M}\varphi^\sharp_{N,M}$ and
$G^\sharp\varphi^\sharp_{N,M}=g_{N,M}\varphi^\sharp_{N,M}$ (on the top slice $\rho=L$ the $\mathsf S^\natural$-terms vanish
and~\eqref{eq:HGsharp} reduces to $q^{-L}h^Q+x(q^{-L}+q^L)$, resp.\ $(q^{-2L}+1)h^Q$).

\begin{proposition}\label{prop:vsharp}
Let $v_\rho:=\sum_ec_{e,\rho}$ be the coefficient sum of $\varphi^\sharp_{N,M}(\rho,\cdot)$ relative to $a_\rho:=x^2q^{2\rho}$; so
$v_L=1$. For $\rho<L$,
\[
v_\rho=\frac{-x\,q^{L+N+M-1}(1-x^2q^{2\rho+2})(1-q^{2\rho+1})}{(1-q^{L-\rho})(1-q^{L+\rho})(1-xq^{N+M-\rho})(1-xq^{N+M+\rho})}\,v_{\rho+1}.
\]
\end{proposition}
\begin{proof}
Write $G_\rho:=\varphi^\sharp(\rho,\cdot)$ and $a:=a_\rho$ (so $a_{\rho+1}=aQ$), and $\mu_j:=v^Q_a((h^Q_a)^jG_\rho)$. The slice $\rho$
of~\eqref{eq:HGsharp} gives
\begin{gather*}
[q^{-\rho}h^Q_a+x(q^{-\rho}+q^\rho)-\varepsilon]G_\rho=xq^\rho(1+xY)G_{\rho+1},\\
(q^{-2\rho}+1)h^Q_aG_\rho-x(1+q^{-1})G_{\rho+1}-x^2(1+q)q^{2\rho}YG_{\rho+1}+qx^2q^{2\rho}YG_{\rho+2}=gG_\rho .
\end{gather*}
Applying $v^Q_a$ (the $Y$-terms have coefficient sum $0$, and $v^Q_a(G_{\rho+1})=(1-aQ)v_{\rho+1}$):
$q^{-\rho}\mu_1+(x(q^{-\rho}+q^\rho)-\varepsilon)\mu_0=xq^\rho(1-aQ)v_{\rho+1}$ and
$(q^{-2\rho}+1)\mu_1-x(1+q^{-1})(1-aQ)v_{\rho+1}=g\mu_0$. Eliminating $\mu_1$ with $c:=q^\rho+q^{-\rho}$ (note
$q^{-\rho}/(q^{-2\rho}+1)=1/c$ and $q^\rho c=1+q^{2\rho}$):
\[\mu_0[g+xc^2-\varepsilon c]=x(1-aQ)(q^{2\rho}-q^{-1})v_{\rho+1}.\]
With $e_n=\sqrt x\,\zeta_n$, $\zeta_n:=z_n+z_n^{-1}$, $z_n:=\sqrt x\,q^n$, the bracket is
$x(c^2-\zeta_N\zeta_Mc+\zeta_N^2+\zeta_M^2-4)=x(c-z_Nz_M-(z_Nz_M)^{-1})(c-z_N/z_M-z_M/z_N)$ (the Fricke relation), i.e.\
$q^{-L-N-M}(1-q^{L-\rho})(1-q^{L+\rho})(1-xq^{N+M-\rho})(1-xq^{N+M+\rho})$.
\end{proof}

\begin{corollary}\label{cor:Vsharp}
$V^\sharp_{N,M}:=v_0$ equals $1$ for $L=0$ and, for $L\ge1$,
\[
V^\sharp_{N,M}=(-x)^Lq^{L(2L-1)+2ML}\,\frac{(1+q^L)\,(x^2Q;Q)_L}{(Q;Q)_L\,(xq^{2M+1};q)_L\,(xq^{2M+L};q)_L}.
\]
\end{corollary}
\begin{proof}
Iterate Proposition~\ref{prop:vsharp}: the product of the factors $(1-q^{2\rho+1})/((1-q^{L-\rho})(1-q^{L+\rho}))$ over
$0\le\rho<L$ is $(q;q^2)_L/((q)_L(q^L;q)_L)=(1+q^L)/(Q;Q)_L$, and the $x$-dependent factors give the rest.
\end{proof}

\begin{lemma}\label{lem:psiu}
$\psi_{N,M}(u):=\sum_{r,s}w(r,s)\varphi^\sharp_{N,M}(r,s)/(q)_r=\varpi_{N,M}V^\sharp_{N,M}/(x^2Q;Q)_\infty$.
\end{lemma}
\begin{proof}
As for Lemma~\ref{lem:omphi}, with the base-$Q$ forms of (B-ii),
\[\textstyle\sum_sw(r,s)\varphi^Q_e(s;a_r)=w(r,e)/(a_rQ;Q)_\infty=w(r,e)(x^2Q;Q)_r/(x^2Q;Q)_\infty,\]
and of (B-i) for $(T^\sharp\varphi^\sharp)(0,\cdot)=\varpi\,\varphi^\sharp(0,\cdot)$, and $v^Q_{x^2}(\varphi^Q_e(\cdot;x^2Q^r))=(x^2Q;Q)_r$.
\end{proof}

\subsection{Biorthogonality, completeness and the spectral expansion}\label{ssec:spec}
Let $\psi_{N,M}(F):=\sum_{r,s}w(r,s)\varphi^\sharp_{N,M}(r,s)F(r,s)$ (over $r\le L$, $s\ge M$; it converges for $F\in\mathcal V$).

\begin{lemma}\label{lem:biorth}
$\psi_a(\varphi_b)=0$ for modes $a\ne b$, and
\[\psi_{N,M}(\varphi_{N,M})=\frac{\varpi_{N,M}}{(x^2q^{2L+2};Q)_{2M}\,(xq^{2M+1};q)_{2L}}.\]
\end{lemma}
\begin{proof}
For $a\ne b$ the sum defining $\psi_a(T\varphi_b)$ is finite ($r\le L_a$, $s\le M_b$), so it can be computed as
$\varpi_b\psi_a(\varphi_b)$ and as $\varpi_a\psi_a(\varphi_b)$, and $\varpi_a\ne\varpi_b$. By the supports
($\varphi_{N,M}$: $R\ge L$, $S\le M$; $\varphi^\sharp_{N,M}$: $R\le L$, $S\ge M$) only $(r,s)=(L,M)$ contributes to
$\psi_{N,M}(\varphi_{N,M})$; there $w=\varpi_{N,M}$, $\varphi_{N,M}(L,M)=\varphi_L(L;xq^{2M})=1/(xq^{2M+1})_{2L}$ and
$\varphi^\sharp_{N,M}(L,M)=1/(x^2q^{2L}Q;Q)_{2M}$.
\end{proof}

\begin{lemma}\label{lem:val}
\begin{enumerate}[label=(\alph*)]
\item $\operatorname{ord}c^{(S)}_e\ge2(M-S)$ for the coefficients of $\varphi_{N,M}$; hence $\operatorname{ord}\varphi_{N,M}(R,S)\ge2(M-S)$.
\item $\operatorname{ord}c_{e,\rho}\ge0$ for $\varphi^\sharp_{N,M}$; hence $\operatorname{ord}\varphi^\sharp_{N,M}(R,S)\ge0$.
\item $\operatorname{ord}\psi_{L'+M',M'}(\Xi_{e,\sigma})\ge e+2\sigma$, with equality and $\psi\ne0$ for $(e,\sigma)=(L',M')$; and
$\psi_{L'+M',M'}(\Xi_{e,\sigma})=0$ unless $e\le L'$ and $\sigma\ge M'$.
\end{enumerate}
\end{lemma}
\begin{proof}
(a) By Lemma~\ref{lem:HGeig} and~\eqref{eq:HGh}, on the slice $S$,
$[\varepsilon_{e+S,S}-\varepsilon_{N,M}]c^{(S)}_e=x^2q^{2S}(xq^{2S+1})_2\cdot$(the $\varphi_e$-coefficient of
$\sum_{e'}c^{(S+1)}_{e'}\tau\varphi_{e'+1}$, relative to $xq^{2S}$). The bracket has order $0$ unless $e+2S=N+M$ (the largest
index on the slice $S$), where it has order $1$ (its $x$-coefficient is $q^{N+M-2S}+q^{2S-N-M}-q^L-q^{-L}\ne0$); that
index is reached only through $\gamma$, and $\operatorname{ord}\gamma\ge1$ ($\gamma$ contains $a=xq^{2S}$). Induction downwards
from $S=M$.
(b) The same with the first equation in the proof of Proposition~\ref{prop:vsharp}: the right side carries a factor
$x$; the bracket $\varepsilon_{\rho+e,e}-\varepsilon_{N,M}$ is nonzero for $\rho<L$ and has order $\le1$ (if its constant
terms $q^{-\rho-2e}$ and $q^{-N-M}$ agree, its $x$-coefficient is $q^\rho+q^{-\rho}-q^L-q^{-L}\ne0$); and the
re-expansion coefficients $c_1,c_2,b_e$ have order $\ge0$.
(c) $\psi(\Xi_{e,\sigma})=\sum_{e\le r\le L'}w(r,\sigma)\varphi^\sharp(r,\sigma)\varphi_e(r;xq^{2\sigma})$ with
$\operatorname{ord}w(r,\sigma)=r+2\sigma$; at $(e,\sigma)=(L',M')$ only $r=L'$ occurs, and the term is the monomial times
unit of Lemma~\ref{lem:biorth}.
\end{proof}

\begin{lemma}[completeness]\label{lem:compl}
If $\Delta\in\mathcal V$ and $\psi_{N,M}(\Delta)=0$ for all $N\ge M\ge0$, then $\Delta=0$.
\end{lemma}
\begin{proof}
Expand $\Delta=\sum\delta_{e,\sigma}\Xi_{e,\sigma}$; the $\delta$'s have order $\ge\inf\operatorname{ord}\Delta$ (Bailey
inversion has $x$-integral coefficients). If $\Delta\ne0$, let $m:=\min\{\operatorname{ord}\delta_{e,\sigma}+e+2\sigma\}$,
attained on a finite nonempty set $Z$. Choose $(e^*,\sigma^*)\in Z$ with $e^*$ minimal and then $\sigma^*$ maximal, and
evaluate $\psi_{e^*+\sigma^*,\sigma^*}(\Delta)=\sum_{e\le e^*,\ \sigma\ge\sigma^*}\psi(\Xi_{e,\sigma})\delta_{e,\sigma}$ (an
$x$-adically convergent sum by Lemma~\ref{lem:val}(c)). The term $(e^*,\sigma^*)$ has order exactly $m$; a term with
$e<e^*$, or $e=e^*$ and $\sigma>\sigma^*$, is not in $Z$ and so has order $\ge\operatorname{ord}\delta_{e,\sigma}+e+2\sigma>m$.
Hence $\psi_{e^*+\sigma^*,\sigma^*}(\Delta)\ne0$.
\end{proof}

\begin{theorem}\label{thm:spec}
In $\mathcal V$ (pointwise, $x$-adically convergent),
\[
u=\sum_{N\ge M\ge0}\tilde u_{N,M}\varphi_{N,M},\qquad
\tilde u_{N,M}:=\frac{\psi_{N,M}(u)}{\psi_{N,M}(\varphi_{N,M})}=\frac{V^\sharp_{N,M}(x^2q^{2L+2};Q)_{2M}(xq^{2M+1};q)_{2L}}{(x^2Q;Q)_\infty},
\]
and for every $k\ge1$
\[
K_k(x)=\sum_{N\ge M\ge0}\varpi_{N,M}^k\,\frac{V_{N,M}V^\sharp_{N,M}(x^2q^{2L+2};Q)_{2M}(xq;q)_{2N}}{(xq;q)_\infty(x^2Q;Q)_\infty}.
\]
\end{theorem}
\begin{proof}
$\operatorname{ord}\tilde u\ge0$ (Corollary~\ref{cor:Vsharp}), so by Lemma~\ref{lem:val}(a) the $(N,M)$-term of
$\tilde U:=\sum\tilde u\varphi$ has order $\ge2(M-S)$ at $(R,S)$, and only $L\le R$, $M\ge S$ contribute; so $\tilde U\in\mathcal V$
converges pointwise. Summability: in the families $(r,s,b)\mapsto w(r,s)\varphi^\sharp_a(r,s)\tilde u_b\varphi_b(r,s)$,
$w(r,s)\tilde u_b\varpi_b^{j}\varphi_b(r,s)/((q)_{R-r}(Q;Q)_{s-S})$ and $w(r,s)\tilde u_b\varpi_b^{j}\varphi_b(r,s)/(Q;Q)_s$ ($j\ge0$), a
nonzero term has $r\ge L_b$, $s\le M_b$ and order $\ge r+2M_b+j(N_b+M_b)$ (Lemma~\ref{lem:val}(a), $\operatorname{ord}\tilde u\ge0$,
$\operatorname{ord}\varphi^\sharp\ge0$), so only finitely many terms lie below any given order. Hence
$\psi_a(\tilde U)=\sum_b\tilde u_b\psi_a(\varphi_b)=\psi_a(u)$ by Lemma~\ref{lem:biorth}, and Lemma~\ref{lem:compl} gives
$u=\tilde U$ (the value of $\tilde u_{N,M}$ by Lemmas~\ref{lem:psiu} and~\ref{lem:biorth}). The same summability gives
$T^{k-1}u=\sum\tilde u\,\varpi^{k-1}\varphi$ (by induction on $k$) and, by~\eqref{eq:Kop} and Lemma~\ref{lem:omphi},
$K_k=\sum\tilde u_{N,M}\varpi^k_{N,M}A_{N,M}/(xq)_\infty$; insert $A=(xq)_{2M}V$ and $(xq)_{2M}(xq^{2M+1})_{2L}=(xq)_{2N}$.
\end{proof}

\subsection{The product identity; proof of Theorem~\ref{thm:F}}\label{ssec:proofF}
Put $t_{N,M}:=\varpi_{N,M}(\beta^B_N\beta^J_M+\beta^B_M\beta^J_N)/(1+\delta_{NM})$. By the computation of $E$ in the proof of (B4),
\begin{multline*}
t_{N,M}=\varpi_{N,M}(-1)^{N+M}q^{N^2+M^2-N-M}\,\frac{(q^N+q^M)(1+xq^{N+M})}{1+\delta_{NM}}\\
\times\frac{(1-xq^{2N})(1-xq^{2M})(x^2;Q)_N(x^2;Q)_M}{(Q;Q)_N(Q;Q)_M(1-x)(1-x^2)}.
\end{multline*}

\begin{lemma}\label{lem:prod}
$V_{N,M}V^\sharp_{N,M}(x^2q^{2L+2};Q)_{2M}(xq;q)_{2N}=t_{N,M}$.
\end{lemma}
\begin{proof}
Collect factors ($N=L+M$). The $Q$-factors: $(Q;Q)_M(Q^{L+1};Q)_M(Q;Q)_L=(Q;Q)_M(Q;Q)_N$. The $x^2$-factors:
$(x^2Q;Q)_L(x^2q^{2L+2};Q)_{2M}/(x^2q^{2N};Q)_M=(x^2Q;Q)_{N-1}(1-x^2q^{2N+2M})$. The remaining $x$-factors combine, via
$(x;q)_M(xq^M;q)_M=(x;q)_{2M}$, $(x;q)_{2M}(1-xq^{2M})(xq^{2M+1};q)_L=(x;q)_{N+M+1}$, $(x^2;Q)_M=(x;q)_M(-x;q)_M$ and
$(x;q)_{2N+1}=(x;q)_{N+M}(1-xq^{N+M})(xq^{N+M+1};q)_{L-1}(1-xq^{2N})$ ($L\ge1$), to the $x$-part of $t_{N,M}$. The monomials
agree ($x^{N+M}(-1)^{N+M}q^{2L^2+4LM+4M^2-L-M}$ on both sides, using $q^N+q^M=q^M(1+q^L)$). For $L=0$, $V^\sharp=1$
matches $(q^N+q^M)/(1+\delta_{NM})=q^M$.
\end{proof}

\begin{proof}[Proof of Theorem~\ref{thm:F}]
By Theorem~\ref{thm:spec} and Lemma~\ref{lem:prod}, and by symmetrisation ($\varpi_{N,M}=\varpi_N\varpi_M$; the double
series converge $x$-adically),
\begin{align*}
(xq)_\infty(x^2Q;Q)_\infty K_k&=\sum_{N\ge M}\varpi_{N,M}^kt_{N,M}=\sum_{n,m\ge0}(\varpi_n\varpi_m)^{k+1}\beta^B_n\beta^J_m\\
&=(xq)_\infty B_{k+1}(x)\cdot(x^2Q;Q)_\infty J_k(x)
\end{align*}
by (B1), (B2), which hold in $\mathbb F$ (§\ref{ssec:transfer}). This holds in $\mathbb F$, hence in $\Z[x][[q]]$.
\end{proof}

\subsection{Identity (iv)}\label{ssec:iv}
Put $O:=PY\mathsf U$, i.e.\ $(OF)(R,S)=q^{R+2S}F(R,S+1)$; $O\mathcal V\subset\mathcal V$.

\begin{lemma}\label{lem:Omop}
$K^\Omega_k(x)=\sum_{i=1}^k\omega(T^{i-1}OT^{k-i}u)$ for $k\ge1$.
\end{lemma}
\begin{proof}
With $F_k:=u$ and $F_i:=TF_{i+1}$, $F_i$ is the function of $(r_i,s_i)$ obtained by summing the levels $i+1,\dots,k$
(their weights and the bonds between the levels $i,\dots,k$); the bond between levels $i$ and $i+1$ is
$1/((q)_{r_i-r_{i+1}}(Q;Q)_{s_{i+1}-s_i})$, and $\omega$ supplies $w(r_1,s_1)/(Q;Q)_{s_1}$. For $i<k$ the $i$-th term of $\Omega$
multiplies the summand by $q^{r_i+2s_i}(1-Q^{s_{i+1}-s_i})$; since $(1-Q^m)/(Q;Q)_m=1/(Q;Q)_{m-1}$ for $m\ge0$ (both sides
vanish for $m=0$, and both bonds vanish for $m<0$), this replaces $F_i=TF_{i+1}$ by $PY\mathsf UTF_{i+1}=OTF_{i+1}$, because
$\sum_{r\le R,\ s\ge S+1}w(r,s)F(r,s)/((q)_{R-r}(Q;Q)_{s-S-1})=(TF)(R,S+1)$. The $k$-th term multiplies $F_k=u$ by
$q^{r_k+2s_k}$, giving $PYu=Ou$ ($u$ does not depend on $S$). For each power of $x$ only finitely many index tuples
contribute, so all rearrangements are finite, coefficient by coefficient.
\end{proof}

Let $h:=h_{xY}=P^{-1}(1-\mathsf S)+xPY$, i.e.\ $(hF)(R,S)=q^{-R}(F(R,S)-F(R-1,S))+xq^{R+2S}F(R,S)$.

\begin{lemma}\label{lem:hcomm}
$Th-hT=-xOT$ on $\mathcal V$.
\end{lemma}
\begin{proof}
By~\eqref{eq:TEW}, $\mathcal W\mathsf S\mathcal W^{-1}=xq^{-1}P^2Y\mathsf S$, so $\mathcal Wh=\tilde h\mathcal W$ with
$\tilde h:=P^{-1}-xq^{-1}PY\mathsf S+xPY$. It suffices to show $hE-E\tilde h=xPY\mathsf UE$; then
$hT-Th=hE\mathcal W-E\tilde h\mathcal W=xPY\mathsf UE\mathcal W=xOT$. The coefficients of $\mathsf S^j\mathsf U^b$, times
$(q)_j(Q;Q)_b$, are: in $hE$, $q^jP^{-1}+xPY$; in $E\tilde h$, $q^jP^{-1}-xq^{-j}Q^b(1-q^j)PY+xq^{-j}Q^bPY=q^jP^{-1}+xQ^bPY$;
in $xPY\mathsf UE$, $xPY(1-Q^b)$. (Restriction to $\N^2$ as in Proposition~\ref{prop:comm}: $h$ reads only $(R,S)$ and
$(R-1,S)$, and $xPY\mathsf U$ reads $E\mathcal WF$ only at $(R,S+1)\in\N^2$.)
\end{proof}

\begin{lemma}\label{lem:omh}
$\omega(hF)=\omega(P^{-1}F)$ for $F\in\mathcal V$.
\end{lemma}
\begin{proof}
Fix $s$. Summation by parts ($F(-1,s)=0$; all sums converge $x$-adically) and $w(r+1,s)=xq^{2r+2s+1}w(r,s)$ give
$\sum_rw(r,s)q^{-r}(F(r,s)-F(r-1,s))=\sum_rw(r,s)F(r,s)(q^{-r}-xq^{r+2s})$. Adding $\sum_rw(r,s)xq^{r+2s}F(r,s)$ gives
$\sum_rw(r,s)q^{-r}F(r,s)$; divide by $(Q;Q)_s$ and sum over $s$.
\end{proof}

\begin{proposition}\label{prop:tele}
For $k\ge1$, $K^\Omega_k(x)=x^{-1}\bigl(K^\flat_k(x)-K_k(x)\bigr)$, where $K^\flat_k:=\omega(P^{-1}T^{k-1}u)$ is $K_k$ with the
extra factor $q^{-r_1}$ in the summand.
\end{proposition}
\begin{proof}
$hu=u+xOu$, because $q^{-R}(1/(q)_R-1/(q)_{R-1})=1/(q)_R$ (also at $R=0$) and $xq^{R+2S}/(q)_R=x(Ou)(R,S)$. For $i<k$,
Lemma~\ref{lem:hcomm} gives $OT^{k-i}u=-x^{-1}(Th-hT)T^{k-i-1}u$. With $a_i:=\omega(T^ihT^{k-1-i}u)$ ($0\le i\le k-1$) the first
$k-1$ terms of Lemma~\ref{lem:Omop} add up to $-x^{-1}\sum_{i=1}^{k-1}(a_i-a_{i-1})=-x^{-1}(a_{k-1}-a_0)$, and
$-x^{-1}a_{k-1}=-x^{-1}\omega(T^{k-1}u)-\omega(T^{k-1}Ou)=-x^{-1}K_k-(\text{the $k$-th term})$. Hence
$K^\Omega_k=x^{-1}(a_0-K_k)$, and $a_0=\omega(hT^{k-1}u)=\omega(P^{-1}T^{k-1}u)$ by Lemma~\ref{lem:omh}. (For $k=1$ the
statement is $\omega(Ou)=x^{-1}(\omega(hu)-\omega(u))$.)
\end{proof}

\begin{lemma}\label{lem:hphi}
$\omega(h\varphi_{N,M})=\dfrac{\varepsilon_{N,M}}{1+x}\,\omega(\varphi_{N,M})$ for all $N\ge M\ge0$.
\end{lemma}
\begin{proof}
On the slice $S=\sigma$, $h$ acts as $h_{xq^{2\sigma}}$ in $R$, so by (B-iii) $h\Xi_{e,\sigma}=(q^{-e}+xq^{2\sigma+e})\Xi_{e,\sigma}$, and
$h\varphi_{N,M}$ is again a finite combination of the $\Xi$'s. By Lemma~\ref{lem:hcomm} and $T\varphi=\varpi\varphi$
($\varpi:=\varpi_{N,M}$), $T(h\varphi)=hT\varphi-xOT\varphi=\varpi(h\varphi-xO\varphi)$, whose slice $S=0$ is
$\varpi(h_xF_0-xq^RF_1(R))$, with $F_0$, $F_1$ as in the proof of Proposition~\ref{prop:Arec}. For $M\ge1$ that proof shows
$v_x(q^RF_1)=0$ and $(1+x)v_x(h_xF_0)=\varepsilon\,v_x(F_0)$. For $M=0$, $F_1=0$, $\varphi_{L,0}=\Xi_{L,0}$,
$h_xF_0=(q^{-L}+xq^L)F_0$ and $\varepsilon_{L,0}=(1+x)(q^{-L}+xq^L)$. In both cases
$v_x((Th\varphi)(\cdot,0))=\varpi\varepsilon A_{N,M}/(1+x)$, and Lemma~\ref{lem:omphi} (applied to $h\varphi$ and to $\varphi$) gives the claim.
\end{proof}

\begin{theorem}\label{thm:specOm}
For $k\ge1$, in $\mathbb F$:
\begin{gather*}
(xq)_\infty(x^2Q;Q)_\infty K^\flat_k=\sum_{N\ge M\ge0}\varpi^k_{N,M}t_{N,M}\frac{\varepsilon_{N,M}}{1+x},\\
(xq)_\infty(x^2Q;Q)_\infty K^\Omega_k=\sum_{N\ge M\ge0}\varpi^k_{N,M}t_{N,M}\frac{\varepsilon_{N,M}-1-x}{x(1+x)}.
\end{gather*}
\end{theorem}
\begin{proof}
By Theorem~\ref{thm:spec}, $T^{k-1}u=\sum_b\tilde u_b\varpi_b^{k-1}\varphi_b$ with the summability stated there, which is not
affected by the factor $q^{-r}$. So $K^\flat_k=\omega(P^{-1}T^{k-1}u)=\sum_b\tilde u_b\varpi_b^{k-1}\omega(P^{-1}\varphi_b)$, and
$\omega(P^{-1}\varphi_b)=\omega(h\varphi_b)=\varepsilon_b\,\omega(\varphi_b)/(1+x)$ by Lemmas~\ref{lem:omh} and~\ref{lem:hphi}. As in the
proof of Theorem~\ref{thm:spec}, $\tilde u_b\varpi_b\,A_b/(xq)_\infty=\varpi_bt_b/((xq)_\infty(x^2Q;Q)_\infty)$ (Lemmas~\ref{lem:omphi}
and~\ref{lem:prod}); this gives the first formula, and the second follows with Proposition~\ref{prop:tele} and
$(xq)_\infty(x^2Q;Q)_\infty K_k=\sum\varpi^kt$.
\end{proof}

\begin{lemma}\label{lem:factor}
For all $N\ge M\ge0$,
$\dfrac{t_{N,M}(\varepsilon_{N,M}-1-x)}{x(1+x)}=\dfrac{\varpi_N\varpi_M\,(b^\Omega_N\beta^J_M+b^\Omega_M\beta^J_N)}{1+\delta_{NM}}$.
\end{lemma}
\begin{proof}
Put $\mu_n:=q^{-2n}+x^2q^{2n}$, so $b^\Omega_n=\beta^B_n(\mu_n+x-1)/(x(1+x))$. Since
$t_{N,M}=\varpi_N\varpi_M(\beta^B_N\beta^J_M+\beta^B_M\beta^J_N)/(1+\delta_{NM})$, the claim is equivalent to
$\beta^B_N\beta^J_M(\varepsilon-\mu_N-2x)+\beta^B_M\beta^J_N(\varepsilon-\mu_M-2x)=0$, $\varepsilon:=\varepsilon_{N,M}$. With $u=q^N$, $v=q^M$,
$\varepsilon=1/(uv)+xu/v+xv/u+x^2uv$, and expanding shows $\varepsilon-\mu_N-2x=(u-v)(1+xu^2)(1-xuv)/(u^2v)$ and
$\varepsilon-\mu_M-2x=(v-u)(1+xv^2)(1-xuv)/(v^2u)$. So the left side is
$\frac{(u-v)(1-xuv)}{uv}\bigl[\beta^B_N(1+xu^2)u^{-1}\beta^J_M-\beta^B_M(1+xv^2)v^{-1}\beta^J_N\bigr]$, and
$\beta^B_n(1+xq^{2n})q^{-n}=(1+x)\beta^J_n$ for all $n\ge0$ (both equal
$(-1)^nq^{n^2-n}(x^2;Q)_n(1-x^2q^{4n})/((Q;Q)_n(1-x))$), so the bracket is $(1+x)(\beta^J_N\beta^J_M-\beta^J_M\beta^J_N)=0$.
\end{proof}

\begin{proof}[Proof of Theorem~\ref{thm:iv}]
By Theorem~\ref{thm:specOm} and Lemma~\ref{lem:factor}, with symmetrisation (the double series converge $x$-adically:
$\operatorname{ord}\varpi_n^{k+1}b^\Omega_n\ge(k+1)n-1$ for $n\ge1$),
\[(xq)_\infty(x^2Q;Q)_\infty K^\Omega_k=\sum_{N\ge M}\frac{(\varpi_N\varpi_M)^{k+1}(b^\Omega_N\beta^J_M+b^\Omega_M\beta^J_N)}{1+\delta_{NM}}
=\Bigl(\sum_n\varpi_n^{k+1}b^\Omega_n\Bigr)\Bigl(\sum_m\varpi_m^{k+1}\beta^J_m\Bigr),\]
and the second factor is $(x^2Q;Q)_\infty J_k(x)$ by (B2). This proves
$(xq)_\infty K^\Omega_k=J_k\sum_n\varpi_n^{k+1}b^\Omega_n$ in $\mathbb F$, hence (§\ref{ssec:transfer}; both sides lie in
$\Z[x][[q]]$ by~\eqref{eq:bOmega}) in $\Z[x][[q]]$, and Proposition~\ref{prop:Fimplies}(b) gives~\eqref{eq:iv}.

It remains to show $(xq;q)_\infty B^{(k)}_{k+1}(x)=\sum_n\varpi_n^{k+1}b^\Omega_n$. Use Bailey pairs relative to $a$ in base $q$,
$\beta_n=\sum_{r\le n}\alpha_r\varphi_r(n;a)$. By the proof of (B1), $\beta_n=1/(Q;Q)_n$,
$\alpha_n=(-1)^nq^{n^2}(a^2;Q)_n(1-aq^{2n})/((Q;Q)_n(1-a))$ is a Bailey pair relative to any $a$; take $a=xq$. Two steps
(L$\infty$) relative to $xq$ give $\alpha''_n=(xq)^{2n}q^{2n^2}\alpha_n$ and
\[\beta''_n=\sum_{n\ge N_k\ge N_{k+1}}\frac{(xq)^{N_k+N_{k+1}}q^{N_k^2+N_{k+1}^2}}{(q)_{n-N_k}(q)_{N_k-N_{k+1}}(Q;Q)_{N_{k+1}}}.\]
By (B-iv) with $c=1$,
$\varphi_d(\cdot;xq)=\frac{1-xq}{1-xq^{2d+1}}\varphi_d(\cdot;x)-\frac{xq^{2d+1}(1-xq)}{1-xq^{2d+1}}\varphi_{d+1}(\cdot;x)$, so
$(\alpha^*,\beta'')$ is a Bailey pair relative to $x$ with
$\alpha^*_e=(1-xq)\alpha''_e/(1-xq^{2e+1})-xq^{2e-1}(1-xq)\alpha''_{e-1}/(1-xq^{2e-1})$, i.e.\ $\alpha^*_0=1$ and, for $e\ge1$,
\begin{align*}
\alpha^*_e&=(-1)^e\Bigl[\frac{x^{2e}q^{3e^2+2e}(x^2Q;Q)_e}{(Q;Q)_e}+\frac{x^{2e-1}q^{3e^2-2e}(x^2Q;Q)_{e-1}}{(Q;Q)_{e-1}}\Bigr]\\
&=\frac{(-1)^ex^{2e-1}q^{3e^2-2e}(x^2Q;Q)_{e-1}(1-y+xy^2-x^3y^3)}{(Q;Q)_e}
\end{align*}
with $y:=q^{2e}$; as $1-y+xy^2-x^3y^3=(1-xy)(1+(x-1)y+x^2y^2)$, \eqref{eq:bOmega} gives $\alpha^*_e=\varpi_e^2b^\Omega_e$ (also for
$e=0$). Now $k-1$ steps (L$\infty$) relative to $x$ give $\alpha^{(k)}_n=\varpi_n^{k-1}\alpha^*_n=\varpi_n^{k+1}b^\Omega_n$, and
$\beta^{(k)}_n$ is the sum defining $B^{(k)}_{k+1}$ with the extra factor $1/(q)_{n-N_1}$ and $N_1\le n$ (for $k=1$ there is
no step and $\beta^{(1)}=\beta''$). Letting $n\to\infty$ as in the proof of Theorem~\ref{thm:B} (for each power of $x$ only
finitely many indices and finitely many $r$ contribute, since $\operatorname{ord}\alpha^{(k)}_r\ge r$) gives
$B^{(k)}_{k+1}(x)/(q)_\infty=\sum_r\alpha^{(k)}_r/((q)_\infty(xq)_\infty)$. With the first part,
$K^\Omega_k=B^{(k)}_{k+1}J_k$.
\end{proof}

At $x=1$ the identity $(q)_\infty B^{(k)}_{k+1}(1)=\vartheta(q^k;q^{2k+4})$ is Bressoud's even-modulus
identity~\cite{Bressoud79,Bressoud80} (quoted in~\cite[(3.6)]{KRTW}); the proof above does not use it.

\subsection{The Hall--Littlewood form of \texorpdfstring{$(F_k)$}{(F\_k)}}\label{ssec:HL}
For $y=(y_1,\dots,y_n)$, Griffin, Ono and Warnaar~\cite[§3.1]{GOW} define
\[
L^{(0)}_m(y;q):=\sum_{r\in\N^n}\frac{\Delta_C(yq^r)}{\Delta_C(y)}\prod_iy_i^{2(m+1)r_i}q^{(m+1)r_i^2+n r_i(r_i-1)/2}
\prod_{i,j}\Bigl(-\frac{y_i}{y_j}\Bigr)^{r_i}\frac{(y_iy_j)_{r_i}}{(qy_i/y_j)_{r_i}},
\]
$\Delta_C(y)=\prod_i(1-y_i^2)\prod_{i<j}(y_i-y_j)(y_iy_j-1)$. For $n=2$, $y=(\sqrt x,-\sqrt x)$ and
$(r_1,r_2)=(n',m')$ one computes
$\Delta_C(yq^r)/\Delta_C(y)=(1-xq^{2n'})(1-xq^{2m'})(q^{n'}+q^{m'})(1+xq^{n'+m'})/(2(1-x)^2(1+x))$ and
$\prod_j(y_1y_j)_{n'}/(qy_1/y_j)_{n'}=(x^2;Q)_{n'}/(Q;Q)_{n'}$, whence
\begin{equation}\label{eq:Lspec}
L^{(0)}_k(\sqrt x,-\sqrt x;q)=L_k(x).
\end{equation}

\emph{The $C_n$ Rogers--Selberg identity} (RS): for $m\ge1$,
\[
\sum_{\lambda:\ \lambda_1\le m}q^{|\lambda|}P'_{2\lambda}(y;q)=\frac{L^{(0)}_m(y;q)}{\prod_i(qy_i^2;q)_\infty\prod_{i<j}(qy_iy_j;q)_\infty},
\]
where $P'_\lambda(y;q)$ is the modified Hall--Littlewood function, i.e.\ $P_\lambda(Y;q)$ for the alphabet
$Y=\{y_iq^j:j\ge0\}$. This follows from the $C_n$ Andrews transformation of Bartlett and
Warnaar~\cite[Theorem~4.2]{BW} and the evaluation of its right-hand side in~\cite[§5.1]{BW} (which uses
their Littlewood-type identity~\cite[Corollary~3.4]{BW}), in the limit $b,c\to\infty$; it is stated
in~\cite[§3.1]{GOW} as the level-$m$ $C_n$ Rogers--Selberg identity, where the infinite products on the
right are omitted from the display. With the products it reduces for $n=1$ to the classical
Rogers--Selberg identity and for $m\to\infty$ to the Littlewood sum; the referees of the first version also
confirmed it numerically at a generic point for $n=2$, $m\le3$ (true with the products, false without them).

\begin{corollary}\label{cor:HL}
\begin{enumerate}[label=(\alph*)]
\item Assume (RS) for $n=2$. Then for every $k\ge1$
\[
(HL_k)\qquad K_k(x)=\sum_{\lambda:\ \lambda_1\le k}q^{|\lambda|}P_{2\lambda}(X_x;q),\qquad X_x=\{\pm\sqrt x\,q^i:i\ge0\}.
\]
\item \textup{(}Corollary~\ref{cor:zq}.\textup{)} Unconditionally, \eqref{eq:ii} for a given $k$ is equivalent to
KRTW's Conjecture~2.2 at $z=q$ for that $k$; so, by Theorem~\ref{thm:k3}, Conjecture~2.2 holds at $z=q$ for every $k$.
\end{enumerate}
\end{corollary}
\begin{proof}
(a) For $y=(\sqrt x,-\sqrt x)$ the product in (RS) is $(xq;q)_\infty^2(-xq;q)_\infty=(xq;q)_\infty(x^2Q;Q)_\infty$,
the prefactor in (B4). By (B4), \eqref{eq:Lspec} and (RS),
$B_{k+1}(x)J_k(x)=\sum_{\lambda_1\le k}q^{|\lambda|}P'_{2\lambda}(\sqrt x,-\sqrt x;q)$, and
$P'_\mu(y_1,y_2;q)=P_\mu(X_x;q)$; Theorem~\ref{thm:F} gives $(HL_k)$.
(b) Griffin, Ono and Warnaar~\cite[Theorem~1.3]{GOW} prove, for positive integers $m,n$ with $n\ge2$ and
$\kappa=2m+2n$,
\[
\sum_{\lambda_1\le m}q^{2|\lambda|}P_{2\lambda}(1,q,q^2,\dots;q^{2n-2})
=\frac{(q^\kappa;q^\kappa)^n_\infty}{(q^2;q^2)_\infty(q)^{n-1}_\infty}\prod_{1\le i<j\le n}\theta(q^{j-i},q^{i+j-1};q^\kappa).
\]
For $n=2$, $m=k$ the right side is the product side of~\eqref{eq:ii}. At $z=q$ the left side of
\eqref{eq:conj22} is the left side above, and the right side of \eqref{eq:conj22} is the left side of
\eqref{eq:ii} (the exponent becomes $(r_i+s_i)^2+s_i^2+r_i+2s_i$).
\end{proof}

So the same $2k$-fold sum carries two Hall--Littlewood interpretations: along KRTW's grading
$z^{\sum(r_i+s_i)}$ (base $q^2$, alphabet $\{q^i\}$; conjectural for general $z$) and along the grading
$x^{\sum(r_i+2s_i)}$ used here (base $q$, alphabet $\{\pm\sqrt xq^i\}$; proved for every $k$, given (RS)). They
meet at $x=z=q$.

\subsection{What remains}\label{ssec:remain}
\begin{enumerate}[leftmargin=2em]
\item KRTW's Conjecture~2.2 for general $z$, their Conjecture~4.9 for $k\ge3$, and Primc's $D_3^{(2)}$ conjecture at
level $k\ge3$ remain open; the identities proved here are their specialisations, not the conjectures themselves.
\item The method of §§\ref{ssec:transfer}--\ref{ssec:iv} uses the transfer operator of the $x$-graded sum. We
have not tried to extend it to KRTW's grading $z^{\sum(r_i+s_i)}$ (Conjecture~2.2 for general $z$), and we do not
know operators commuting with the corresponding transfer operator.
\item Identities \eqref{eq:ii}--\eqref{eq:iv} are not formalized in Lean (§\ref{sec:lean}). Beyond the Bailey-pair
machinery already formalized for~\eqref{eq:i}, a formalization would need the $x$-adic convergence and
completeness arguments of §\ref{ssec:spec}.
\end{enumerate}

\section{The Lean formalization}\label{sec:lean}

\subsection{What is formalized}
The Lean~4 project in the public repository at \repo\fpath{lean/lean-krtw/} proves, for every $k$, the printed identity~\eqref{eq:i} in product
form, and the first equality of Theorem~\ref{thm:x} for an arbitrary field $K$ and an arbitrary power
series $x\in K[[q]]$. All series are elements of Mathlib's \texttt{PowerSeries} (the variable \texttt{X}
plays the role of~$q$). Pins: Lean \texttt{leanprover/lean4:v4.33.0-rc2}, Mathlib commit
\texttt{51e6992efd06\allowbreak126df61a496bebf8\allowbreak f49482a4e129}. The library has about 2900 lines. The proofs use
no \texttt{sorry}, no \texttt{axiom} and no \texttt{native\_decide}; for each of the nine statements below
\verb|#print axioms| gives \texttt{[propext, Classical.choice, Quot.sound]}.

\subsection{The definitions}
The trusted part is the file \fpath{KRTWStatement.lean}, which imports only Mathlib and contains eleven
definitions and the nine statements (with proofs replaced by \texttt{sorry} on purpose). The
definitions are (verbatim up to docstrings; the comments are ours):
\begin{lstlisting}
variable {K : Type*} [Field K]
noncomputable def qPoch (n : ℕ) : K⟦X⟧ := ∏ i ∈ range n, (1 - X ^ (i + 1))         -- (q;q)_n
noncomputable def q2Poch (n : ℕ) : K⟦X⟧ := ∏ i ∈ range n, (1 - X ^ (2 * (i + 1)))  -- (q^2;q^2)_n
def nxt {k : ℕ} (r : Fin k → ℕ) (i : Fin k) : ℕ :=                                  -- r_{i+1}, r_{k+1} := 0
  if h : (i : ℕ) + 1 < k then r ⟨i + 1, h⟩ else 0
def prv {k : ℕ} (s : Fin k → ℕ) (i : Fin k) : ℕ :=                                  -- s_{i-1}, s_0 := 0
  if h : (i : ℕ) = 0 then 0 else s ⟨i - 1, by omega⟩
def box (k N : ℕ) : Finset (Fin k → ℕ) := Fintype.piFinset fun _ => range (N + 1)
noncomputable def krtwTerm (x : K⟦X⟧) {k : ℕ} (r s : Fin k → ℕ) : K⟦X⟧ :=
  ∏ i : Fin k, x ^ (r i + 2 * s i) * X ^ ((r i + s i) ^ 2 + s i ^ 2 + s i)
    * (qPoch (r i - nxt r i))⁻¹ * (q2Poch (s i - prv s i))⁻¹
noncomputable def krtwSum (x : K⟦X⟧) (k : ℕ) : K⟦X⟧ :=
  PowerSeries.mk fun N => ∑ r ∈ box k N, ∑ s ∈ box k N,
    if (∀ i, nxt r i ≤ r i) ∧ (∀ i, prv s i ≤ s i) then coeff N (krtwTerm x r s) else 0
noncomputable def melzerTerm (x : K⟦X⟧) {k : ℕ} (M : Fin k → ℕ) : K⟦X⟧ :=
  ∏ i : Fin k, x ^ (M i) * X ^ (M i * (M i + 1) / 2) * (qPoch (M i - nxt M i))⁻¹
noncomputable def melzerSum (x : K⟦X⟧) (k : ℕ) : K⟦X⟧ :=
  PowerSeries.mk fun N => ∑ M ∈ box k N,
    if ∀ i, nxt M i ≤ M i then coeff N (melzerTerm x M) else 0
noncomputable def qPochInf (a d : ℕ) : K⟦X⟧ :=                                   -- (q^a;q^d)_∞
  PowerSeries.mk fun N => coeff N (∏ i ∈ range (N + 1), (1 - X ^ (a + d * i) : K⟦X⟧))
noncomputable def thetaSum (k : ℕ) : K⟦X⟧ :=
  PowerSeries.mk fun N => ∑ n ∈ range (N + 1),
    coeff N ((-1 : K⟦X⟧) ^ n * X ^ (((k + 2) * n ^ 2 + k * n) / 2) * (1 - X ^ (2 * n + 1)))
\end{lstlisting}

\emph{How the definitions match the printed identity.} A tuple $(r_1,\dots,r_k)$ is a function
\lstinline{r : Fin k → ℕ} with \texttt{r 0} $=r_1$; \texttt{nxt r i} is $r_{i+1}$ with $r_{k+1}=0$ and
\texttt{prv s i} is $s_{i-1}$ with $s_0=0$. \texttt{krtwTerm x r s} is the summand of the left side of
\eqref{eq:xref}, so at $x=1$ it is the summand of~\eqref{eq:i}. The symbol \lstinline{⁻¹} is the inverse in
$K[[q]]$ of a series with constant term~$1$. Each side is defined \emph{by its coefficients as a finite
sum}: the coefficient of $q^N$ of \texttt{krtwSum} is the sum, over the tuples with all entries $\le N$
that satisfy the chain conditions $r_{i+1}\le r_i$ and $s_{i-1}\le s_i$, of the coefficient of $q^N$ of
the summand. Natural-number subtraction is truncated in Lean, but it is only evaluated under the chain
conditions, where it is the true difference. The bound $N$ loses nothing: the theorems
\texttt{coeff\_krtwTerm\_eq\_zero} and \texttt{coeff\_melzerTerm\_eq\_zero} below say that a summand
with an entry $>N$ has zero coefficient of $q^N$ (because $(r_i+s_i)^2+s_i^2+s_i\ge\max(r_i,s_i)$ and
$M_i(M_i+1)/2\ge M_i$). KRTW print the sum over all $r_i,s_i\ge0$ with $1/(q)_n=0$ for $n<0$; this is the
same sum. \texttt{qPochInf a d} is $(q^a;q^d)_\infty$ for $a,d\ge1$: its coefficient of $q^N$ is that of the
finite product over $i\le N$, and \texttt{coeff\_qPochInf} says that every longer finite product has the
same coefficient of $q^N$. \texttt{thetaSum k} is $\sum_{n\in\Z}(-1)^nq^{((k+2)n^2+kn)/2}$ with the terms
$n\ge0$ and $-n-1$ paired by hand (the exponent at $-n-1$ is that at $n$ plus $2n+1$, and
$(k+2)n^2+kn$ is even).

\subsection{The theorems}
In \fpath{KRTWSolution.lean} (namespace \texttt{KRTW}) the nine statements are proved (shown verbatim up to docstrings; comments ours):
\begin{lstlisting}
-- Theorem 1.1, the printed identity (i), for every k:
theorem krtw_identity (k : ℕ) :
    krtwSum (1 : ℚ⟦X⟧) k
      = qPochInf 1 (k + 2) * qPochInf (k + 1) (k + 2) * qPochInf (k + 2) (k + 2)
        * (qPochInf 1 2 * qPochInf 1 1)⁻¹
-- the first equality of Theorem 1.2, for an arbitrary field and an arbitrary series x:
theorem krtw_sum_eq_melzer_sum_refined {K : Type*} [Field K] (x : K⟦X⟧) (k : ℕ) :
    krtwSum x k = melzerSum x k
theorem krtw_sum_eq_melzer_sum (k : ℕ) :
    krtwSum (1 : ℚ⟦X⟧) k = melzerSum (1 : ℚ⟦X⟧) k
-- output of the Bailey chain, the triple-product instance, the known k-fold evaluation:
theorem krtw_identity_theta (k : ℕ) :
    krtwSum (1 : ℚ⟦X⟧) k * (qPochInf 1 2 * qPochInf 1 1) = thetaSum k
theorem theta_eq_triple_product (k : ℕ) :
    (thetaSum k : ℚ⟦X⟧) = qPochInf 1 (k + 2) * qPochInf (k + 1) (k + 2) * qPochInf (k + 2) (k + 2)
theorem melzer_identity (k : ℕ) :
    melzerSum (1 : ℚ⟦X⟧) k
      = qPochInf 1 (k + 2) * qPochInf (k + 1) (k + 2) * qPochInf (k + 2) (k + 2)
        * (qPochInf 1 2 * qPochInf 1 1)⁻¹
-- the truncations in the definitions are harmless:
theorem coeff_krtwTerm_eq_zero {K : Type*} [Field K] (x : K⟦X⟧) {k : ℕ} (r s : Fin k → ℕ) (N : ℕ)
    (h : r ∉ box k N ∨ s ∉ box k N) : coeff N (krtwTerm x r s) = 0
theorem coeff_melzerTerm_eq_zero {K : Type*} [Field K] (x : K⟦X⟧) {k : ℕ} (M : Fin k → ℕ) (N : ℕ)
    (h : M ∉ box k N) : coeff N (melzerTerm x M) = 0
theorem coeff_qPochInf {K : Type*} [Field K] (a d : ℕ) (ha : 1 ≤ a) (hd : 1 ≤ d) (N M : ℕ)
    (h : N ≤ M) : coeff N (qPochInf a d : K⟦X⟧) = coeff N (∏ i ∈ range (M + 1), (1 - X ^ (a + d * i) : K⟦X⟧))
\end{lstlisting}
The right side of \texttt{krtw\_identity} is
$(q;q^{k+2})_\infty(q^{k+1};q^{k+2})_\infty(q^{k+2};q^{k+2})_\infty\cdot((q;q^2)_\infty(q;q)_\infty)^{-1}$,
the product side of~\eqref{eq:i}. The identity is stated over~$\Q$; since every coefficient is a
rational number determined by finitely many integer operations, this is the identity of~\cite{KRTW} in
$\Z[[q]]$. Likewise, \texttt{krtw\_sum\_eq\_melzer\_sum\_refined} with $K=\Q(x)$ and the constant series $x$ gives
the first equality of Theorem~\ref{thm:x} in $\Z[x][[q]]$, because $\Z[x][[q]]\to\Q(x)[[q]]$ is injective.
(The solution file also contains auxiliary lemmas \texttt{krtw\_product\_field}, \texttt{krtw\_theta\_field}
proving the product statements over an arbitrary field; they are not among the nine audited statements.)

\subsection{How the proof is organised, and how it differs from the text}
The sums are formed in $K[[q]]$ with the coefficientwise topology (Mathlib's \texttt{tsum} for the
product topology of the discrete topology on~$K$); $x$ is an arbitrary power series, so the
substitutions $x\mapsto xq^c$ are evaluations of a function at another point. Theorem~\ref{thm:R}, the
limit $R\to\infty$, the regrouping~\eqref{eq:31} and the exponent bookkeeping~\eqref{eq:01} are
formalized as in the text (files \fpath{KRTW/Theorem1.lean}, \fpath{KRTW/Limit.lean}). Lemma~\ref{lem:key}
is formalized in its two forms \eqref{eq:star2}, \eqref{eq:starinf} by a different route
(\fpath{KRTW/Lemma4.lean}): with $c(R,n)=\sum_j(-1)^jq^{Rj}/((q)_j(q)_{n-j})$ one has
$c(R+1,n+1)=c(R,n+1)+q^Rc(R,n)$ and $(1-q^{n+1})c(R,n+1)=c(R+1,n)-q^Rc(R,n)$, hence $c(0,2b)=1/(q^2;q^2)_b$,
$c(0,2b+1)=0$; both sides of \eqref{eq:star2} and \eqref{eq:starinf} satisfy
$F(R+1;x)=F(R;x)+xq^{R+1}F(R;xq)$, which reduces them to $R=0$, where a recurrence in $M$ finishes the
proof. Thus Lemmas~\ref{lem:L2} (for $R>0$) and~\ref{lem:L3} and the finite form \eqref{eq:star} are not
formalized as stated; their consequences \eqref{eq:star2}, \eqref{eq:starinf} are. For the $k$-fold sum
(\fpath{KRTW/Bailey.lean}, \fpath{KRTW/JTP.lean}, \fpath{KRTW/Stage2.lean}) the Bailey pair relation is
multiplied through by $(aq;q)_{2n}$ so that no inverse of a Pochhammer symbol with a generic argument is
needed; Bailey's lemma is derived from the $q$-Chu--Vandermonde sum, which is proved; the chain is
formalized for arbitrary~$x$ but the limit $n\to\infty$ only at $x=1$; and the Jacobi triple product is
proved, from the $q$-binomial theorem, only in the instance $(q,q^{k+1},q^{k+2};q^{k+2})_\infty$ (Mathlib
has no triple product). No gap or error was found in the text in the parts that were formalized.

\subsection{What is not formalized}
The $x$-refined closed form \eqref{eq:melclosed} (only its value at $x=1$ is formalized);
Corollaries~\ref{cor:z1}, \ref{cor:zq}, \ref{cor:y}, \ref{cor:HL}; all of §\ref{sec:k3}; novelty
claims. Whether the eleven definitions transcribe the printed display of~\cite{KRTW} is a matter of
reading them (above); a computable mirror (\fpath{scripts/Sanity.lean}, which calls the very
\texttt{KRTW.nxt}, \texttt{KRTW.prv} of the statement file) checks with \texttt{\#guard} that for
$k=1,2,3$ the first $31$ coefficients of \texttt{krtwSum}, \texttt{melzerSum} and of the product equal
integers computed independently by a plain Python enumeration of the printed sums
(\fpath{scripts/sanity_independent.py}); for $k=3$ they begin $1,1,2,4,5,8,12,17,24,34,\dots$. It also
checks that \texttt{krtwSum} times $(q;q^2)_\infty(q;q)_\infty$ equals \texttt{thetaSum} for $k\le3$, and that
\texttt{thetaSum} equals the triple product for $k=3$. Its two negative controls (a product with the wrong
modulus, and \texttt{krtwSum} for $k=2$ against \texttt{melzerSum} for $k=3$) are coarse: they would not
detect, e.g., a wrong exponent in a definition; that is what reading the definitions is for. (The mirror is linked to the noncomputable definitions by
inspection, not by a proof.)

\subsection{How to rebuild and check}
{\footnotesize
\begin{verbatim}
curl -sSfL https://raw.githubusercontent.com/leanprover/elan/master/elan-init.sh \
    | sh -s -- -y --default-toolchain none
export PATH="$HOME/.elan/bin:$PATH"
cd 2026-11-number-theory-q-series/lean/lean-krtw
lake exe cache get    # Mathlib binaries
lake build            # 'declaration uses sorry' only from KRTWStatement.lean
lake env lean scripts/PrintAxioms.lean
./scripts/compare_statements.sh
lake env lean scripts/Sanity.lean
python3 scripts/sanity_independent.py
\end{verbatim}}
The outputs are saved in \fpath{logs/}. The script \fpath{compare_statements.sh} checks that
\fpath{KRTWStatement.lean} imports only Mathlib, that its definition block coincides byte for byte with
\fpath{KRTW/Defs.lean}, that each statement occurs byte for byte in \fpath{KRTWSolution.lean}, that the
solution side contains no \texttt{sorry}, \texttt{axiom} or \texttt{native\_decide}, and that
\texttt{\#check} and \texttt{\#print} give the same output from both modules. (A configuration file
\fpath{comparator.json} for the stricter tool \texttt{leanprover/comparator} is included; that tool was not
run.) A reader who accepts the Lean kernel and the eleven definitions above therefore obtains
Theorem~\ref{thm:main} independently of §§\ref{sec:setting}--\ref{sec:x}.

\section{Numerical verification}\label{sec:num}
None of the following is used in the proofs. Paths are relative to \repo.
\begin{itemize}[leftmargin=1.5em]
\item \fpath{02-exceptional-nahm-sums/code/krtw_general_certify.py} (log
\fpath{02-exceptional-nahm-sums/certificates/krtw_general_certify.txt}), exact arithmetic: Lemma~\ref{lem:L2} for
$R\le8$, $L\le10$; Lemma~\ref{lem:L3} for $R,d\le7$; \eqref{eq:star} as a polynomial identity in
$\Z[x,q]$ for $R,M\le7$; \eqref{eq:star2}, \eqref{eq:starinf} for $R,M\le8$ mod $q^{36}$;
Theorem~\ref{thm:R} for $j\le4$, $R\le6,5,4,3$; Theorem~\ref{thm:x} for $k\le7$, all powers of $x$, mod
$q^{36}$; Corollary~\ref{cor:y} for $k\le3$ at six values of $y$; the Bailey pairs of
Proposition~\ref{prop:bailey} for $i\le3$, $n\le4$; \eqref{eq:melclosed} for $k\le7$; $\Mel_k(1)$ equal to the product for $k\le20$ and \eqref{eq:i} for $k\le8$, mod $q^{300}$;
Remark~\ref{rem:X1} for $k\le6$ mod $q^{80}$.
\item \fpath{02-exceptional-nahm-sums/code/krtw_four_test.py} (log \fpath{certificates/krtw_four_test.txt}):
all four identities \eqref{eq:i}--\eqref{eq:iv} hold modulo $q^{400}$ for $1\le k\le8$ (the $2k$-fold sums
are computed by a level-by-level recursion, compared with brute-force enumeration for $k\le3$ mod $q^{30}$).
\item \fpath{02-exceptional-nahm-sums/code/krtw_Fk_test.py} (log \fpath{certificates/krtw_Fk_test.txt}):
$(F_k)$ and the identity $K^\Omega_k=B^{(k)}_{k+1}J_k$ of Theorem~\ref{thm:iv} for $k\le7$, all powers of $x$, mod $q^{44}$;
\fpath{02-exceptional-nahm-sums/code/krtw_struct_certify.py} (log
\fpath{certificates/krtw_struct_certify.txt}): Theorem~\ref{thm:B} and $(F_k)$ for $k\le7$, all powers of
$x$, mod $q^{90}$.
\item \fpath{02-exceptional-nahm-sums/code/krtw_moves_certify.py}: the move chains of the first version
for $(F_2)$, $(F_3)$, with the exact precondition of every DEFUSE, the matching of the end data, and an
independent series evaluation of all four ends ($x^0,\dots,x^{16}$ mod $q^{60}$).
\item \fpath{01-krtw-identities/all-k/code/certify_F_all_k.py} (log \fpath{01-krtw-identities/all-k/certificates/certify_F_all_k.txt}),
exact arithmetic: Proposition~\ref{prop:comm} and the transposed relations~\eqref{eq:HGsharp} on point masses at
two rational points; (B-v) symbolically; the eigen-relations of Lemma~\ref{lem:HGeig} and of $H'^\sharp$, $G^\sharp$ on
eigenvectors computed independently from the triangular systems of $T$ and $T^\sharp$; Corollaries~\ref{cor:V}
and~\ref{cor:Vsharp} against those eigenvectors ($N\le7$); Propositions~\ref{prop:Arec} and~\ref{prop:vsharp}
($N\le8$); Lemma~\ref{lem:biorth} (all pairs with $N\le3$); Lemma~\ref{lem:prod} ($N\le8$); Theorem~\ref{thm:F} for
$k\le6$, all powers of $x$, mod $q^{30}$; the expansion of $u$ in Theorem~\ref{thm:spec}. Negative controls are
detected.
\item \fpath{01-krtw-identities/all-k/code/certify_iv_all_k.py} (log \fpath{01-krtw-identities/all-k/certificates/certify_iv_all_k.txt}):
Lemmas~\ref{lem:hcomm} and~\ref{lem:omh} on point masses (exact); Proposition~\ref{prop:tele} for $k\le8$, all powers
of $x$, mod $q^{40}$, with $\Omega$ expanded literally; Lemma~\ref{lem:hphi} for $N\le5$; Lemma~\ref{lem:factor} for
$N\le12$ and $\alpha^*_n=\varpi_n^2b^\Omega_n$ for $n\le12$ (symbolic); Theorem~\ref{thm:specOm} and the $x$-refinement of
Theorem~\ref{thm:iv} for $k\le8$, all powers of $x$, mod $q^{40}$; the triple-product evaluation in
Proposition~\ref{prop:Fimplies}(b) for $k\le10$ mod $q^{200}$. Negative controls are detected.
\item Separate code, \fpath{01-krtw-identities/all-k/code/series_check_q300.py}, with log
\fpath{01-krtw-identities/all-k/certificates/series_check_q300.txt}: \eqref{eq:ii}, \eqref{eq:iii}, \eqref{eq:iv} in KRTW's
printed form, with $\Omega$ expanded literally, mod $q^{300}$ for $k\le8$.
\item Independent referee computations (reports \fpath{01-krtw-identities/REFEREE-*.md}; code written
from the statements; the code of one referee is archived in \fpath{02-exceptional-nahm-sums/code/krtw_referee/},
e.g.\ \fpath{chains.py}, \fpath{hl_direct.py}, \fpath{rs_generic.py}; the code of the other three was not
archived): \eqref{eq:i} by direct enumeration for $k\le4$ mod $q^{500}$ and for $k\le8$ to
orders between $q^{400}$ and $q^{140}$; Theorem~\ref{thm:R} for $j\le5$, $R\le8$ mod $q^{60}$;
Lemma~\ref{lem:key} exactly for $R,M\le8$; the Hall--Littlewood side of Corollary~\ref{cor:z1} for
$k\le4$ mod $q^{26}$; all $30$ intermediate data of the chains ($x^0,\dots,x^{16}$ mod $q^{100}$);
\eqref{eq:ii}, \eqref{eq:iii} for $k=3$ in KRTW's printed form mod $q^{150}$; $(HL_k)$ for $k\le3$ directly
from Kostka--Foulkes polynomials.
\item Independent referee computations for §\ref{sec:k3} (reports \fpath{01-krtw-identities/all-k/REFEREE-*.md}, code
archived in \fpath{01-krtw-identities/all-k/referee*/}): \eqref{eq:ii}, \eqref{eq:iii} mod $q^{300}$ for $k\le12$ and
\eqref{eq:iv} mod $q^{401}$ for $k\le14$, from the printed sums; Theorem~\ref{thm:F}, Proposition~\ref{prop:tele},
Theorem~\ref{thm:specOm} and Theorem~\ref{thm:iv} as series in $x$ and $q$ (mod $q^{81}$ for $k\le8$, mod
$q^{61}$ for $k\le12$); the commutation relations on point masses at three rational points; eigenvectors of $T$
and $T^\sharp$ computed pointwise in position space, with all the eigen-relations, values at infinity and
pairings of §§\ref{ssec:basis}--\ref{ssec:spec} and Lemma~\ref{lem:hphi}; the sharpness of the valuation bound of
Lemma~\ref{lem:val}(a); the expansion of $u$ (exactly modulo $x^{14}$); Lemma~\ref{lem:prod} and Lemma~\ref{lem:factor}
symbolically. The second read of §\ref{sec:k3} also found \eqref{eq:ii}--\eqref{eq:iv} mod $q^{400}$ for $k\le10$, and
Theorems~\ref{thm:F}, \ref{thm:iv} as series in $x$ and $q$ mod $q^{100}$ for $k\le8$.
\end{itemize}

\emph{Concordance.} The certificates, the referee reports and the Lean docstrings use the numbering of the
source files \fpath{01-krtw-identities/KRTW-general.md} and \fpath{KRTW-F-moves.md}: there, Lemma~4 is
Lemma~\ref{lem:key}, Theorem~1 is Theorem~\ref{thm:R}, Theorem~2 is Theorem~\ref{thm:x}, Theorem~3 is
Theorem~\ref{thm:main}, Proposition~6 is Proposition~\ref{prop:bailey}, Corollary~4 is
Corollary~\ref{cor:z1}, Corollary~5 is Corollary~\ref{cor:y}; Theorem~B is Theorem~\ref{thm:B}, Corollary~C is
Corollary~\ref{cor:HL} (their Theorem~A, the case $k\le3$ of Theorem~\ref{thm:F}, is superseded). The files
\fpath{01-krtw-identities/all-k/proofs/F-all-k.md} and \fpath{iv-all-k.md}, with their certificates and reports, write $\lambda$
for $\varpi$, $W$ for $\mathcal W$, $W_{N,M}$ for $A_{N,M}$ and $\mathfrak u$ for $\mathsf U$, and number as follows. In
\fpath{F-all-k.md}: Proposition~2.1 is Proposition~\ref{prop:comm}, Proposition~3.1 is Proposition~\ref{prop:eig},
Lemma~3.2 is Lemma~\ref{lem:HGeig}, Proposition~4.1 is Proposition~\ref{prop:Arec}, Corollary~4.2 is
Corollary~\ref{cor:V}, Lemma~4.3 is Lemma~\ref{lem:omphi}, Proposition~5.1 is Proposition~\ref{prop:vsharp},
Corollary~5.2 is Corollary~\ref{cor:Vsharp}, Lemma~5.3 is Lemma~\ref{lem:psiu}, Lemmas~6.1--6.3 are
Lemmas~\ref{lem:biorth}--\ref{lem:compl}, Theorem~6.4 is Theorem~\ref{thm:spec}, Lemma~7.1 is Lemma~\ref{lem:prod},
Theorems~1 and~2 are Theorems~\ref{thm:F} and~\ref{thm:k3}. In \fpath{iv-all-k.md}: Lemma~1.1 is
Lemma~\ref{lem:Omop}, Lemma~2.1 is Lemma~\ref{lem:hcomm}, Lemma~2.2 is Lemma~\ref{lem:omh}, Proposition~2.3 is
Proposition~\ref{prop:tele}, Lemma~3.1 is Lemma~\ref{lem:hphi}, Theorem~3.2 is Theorem~\ref{thm:specOm}, Lemma~4.1 is
Lemma~\ref{lem:factor}, and Theorems~4--6 are Theorem~\ref{thm:iv} together with Proposition~\ref{prop:Fimplies}(b).

\section*{Acknowledgements}
The results, proofs, computer code and Lean formalization in this paper were produced by AI (Anthropic's
Claude models) directed by the author. Their correctness was checked by independent AI referees
(separate instances instructed to find errors, each writing its own code; seven reports, included in
the repository; two independent reads of each proof, §\ref{ssec:reads}) and, for the formalized statements
(identity~\eqref{eq:i}), by the Lean kernel. No human expert has refereed
this work yet; corrections are welcome. We thank the authors of~\cite{KRTW} for stating the identities
precisely, and the developers of Lean and Mathlib.

\begin{thebibliography}{99}

\bibitem{AAB} A.~K. Agarwal, G.~E. Andrews and D.~M. Bressoud, \emph{The Bailey lattice}, J. Indian Math.
Soc. (N.S.) \textbf{51} (1987), 57--73.

\bibitem{Andrews1984} G.~E. Andrews, \emph{Multiple series Rogers--Ramanujan type identities}, Pacific
J. Math. \textbf{114} (1984), 267--283.

\bibitem{AndrewsBook} G.~E. Andrews, \emph{The Theory of Partitions}, Encyclopedia of Mathematics and its
Applications~2, Addison-Wesley, 1976; reprinted Cambridge University Press, 1998.

\bibitem{Andrews2010} G.~E. Andrews, \emph{Parity in partition identities}, Ramanujan J. \textbf{23}
(2010), 45--90.

\bibitem{Bailey} W.~N. Bailey, \emph{Identities of the Rogers--Ramanujan type}, Proc. London Math. Soc.
(2) \textbf{50} (1949), 1--10.

\bibitem{BW} N.~Bartlett and S.~O. Warnaar, \emph{Hall--Littlewood polynomials and characters of affine
Lie algebras}, Adv. Math. \textbf{285} (2015), 1066--1105; arXiv:1304.1602.

\bibitem{Bartholdi} L.~Bartholdi, \emph{Automatic partition identities}, arXiv:2611.04020 (2026).

\bibitem{Bressoud79} D.~M. Bressoud, \emph{A generalization of the Rogers--Ramanujan identities for all
moduli}, J. Combin. Theory Ser. A \textbf{27} (1979), 64--68.

\bibitem{Bressoud80} D.~M. Bressoud, \emph{Analytic and combinatorial generalizations of the
Rogers--Ramanujan identities}, Mem. Amer. Math. Soc. \textbf{24} (1980), no.~227.

\bibitem{BIS} D.~Bressoud, M.~Ismail and D.~Stanton, \emph{Change of base in Bailey pairs}, Ramanujan J.
\textbf{4} (2000), 435--453; arXiv:math/9909053.

\bibitem{CW} Z.~Cao and L.~Wang, \emph{Some new modular Nahm sums of ranks 3 and 4}, arXiv:2606.13590
(2026); see also arXiv:2412.15767 and arXiv:2508.12468.

\bibitem{CLSXY} S.~Chern, Z.~Li, D.~Stanton, T.~Xue and A.~J. Yee, \emph{The Ariki--Koike algebras and
Rogers--Ramanujan type partitions}, J. Algebraic Combin. \textbf{60} (2024), no.~2, 491--540;
arXiv:2209.07713.

\bibitem{Lean4} L.~de~Moura and S.~Ullrich, \emph{The Lean~4 theorem prover and programming language},
in: Automated Deduction -- CADE 28, Lecture Notes in Comput. Sci. \textbf{12699}, Springer, 2021,
625--635.

\bibitem{Ding} X.~Ding, \emph{A Lie-theoretic approach to the type $C_n^{(1)}$ partitions of
Capparelli, Meurman, Primc and Primc}, arXiv:2611.10654 (2026).

\bibitem{DJ} J.~Dousse and J.~Jang, \emph{Andrews--Gordon type identities with parity restrictions
through particle motion}, arXiv:2603.05300 (2026).

\bibitem{EKL} T.~Ekholm, P.~Kucharski and P.~Longhi, \emph{Multi-cover skeins, quivers, and 3d
$\mathcal N=2$ dualities}, J. High Energy Phys. 2020, no.~2, 018; arXiv:1910.06193.

\bibitem{GR} G.~Gasper and M.~Rahman, \emph{Basic Hypergeometric Series}, 2nd ed., Encyclopedia of
Mathematics and its Applications~96, Cambridge University Press, 2004.

\bibitem{GOW} M.~J. Griffin, K.~Ono and S.~O. Warnaar, \emph{A framework of Rogers--Ramanujan identities
and their arithmetic properties}, Duke Math. J. \textbf{165} (2016), no.~8, 1475--1527; arXiv:1401.7718.

\bibitem{KL} S.~Kanade and J.~Lovejoy, \emph{Bailey pairs and an identity of
Chern--Li--Stanton--Xue--Yee}, SIGMA \textbf{21} (2025), 021; arXiv:2410.15408.

\bibitem{KR} S.~Kanade and M.~C. Russell, \emph{Tight cylindric partitions}, arXiv:2508.15113 (2025).

\bibitem{KRTW} S.~Kanade, M.~C. Russell, S.~Tsuchioka and S.~O. Warnaar, \emph{Remarks on the
conjectures of Capparelli, Meurman, Primc and Primc}, Selecta Math. (N.S.) \textbf{32} (2026), no.~5,
Paper No.~100; arXiv:2404.03851.

\bibitem{KimYee} S.~Kim and A.~J. Yee, \emph{Rogers--Ramanujan--Gordon identities, generalized
G\"ollnitz--Gordon identities, and parity questions}, J. Combin. Theory Ser. A \textbf{120} (2013).

\bibitem{KLNS} P.~Kucharski, H.~Larragu\'ivel, D.~Noshchenko and P.~Su\l kowski, \emph{Unlinking symmetric
quivers}, arXiv:2312.14905 (2023).

\bibitem{Kursungoz} K.~Kur\c{s}ung\"oz, \emph{Combinatorial construction of known positive series for
partition classes defined by Capparelli, Meurman, Primc, and Primc in the $k=1$ case},
arXiv:2605.13096 (2026).

\bibitem{Macdonald} I.~G. Macdonald, \emph{Symmetric Functions and Hall Polynomials}, 2nd ed., Oxford
University Press, 1995.

\bibitem{Mathlib} The mathlib Community, \emph{The Lean mathematical library}, in: Proceedings of the 9th
ACM SIGPLAN International Conference on Certified Programs and Proofs (CPP 2020), ACM, 2020, 367--381.

\bibitem{Melzer} E.~Melzer, \emph{Supersymmetric analogs of the Gordon--Andrews identities, and related
TBA systems}, arXiv:hep-th/9412154 (1994).

\bibitem{MP} A.~Meurman and M.~Primc, \emph{Annihilating fields of standard modules of
$\widetilde{\mathfrak{sl}}(2,\mathbb C)$ and combinatorial identities}, Mem. Amer. Math. Soc.
\textbf{137} (1999), no.~652.

\bibitem{CRW} Z.~Cao, H.~Rosengren and L.~Wang, \emph{On some double Nahm sums of Zagier}, J. Combin.
Theory Ser. A \textbf{202} (2024), Paper No.~105819; arXiv:2303.01333.

\bibitem{Russell} M.~C. Russell, \emph{Companions to the Andrews--Gordon and Andrews--Bressoud identities
and recent conjectures of Capparelli, Meurman, Primc, and Primc}, SIGMA \textbf{22} (2026), 046;
arXiv:2306.16251.

\bibitem{Slater} L.~J. Slater, \emph{Further identities of the Rogers--Ramanujan type}, Proc. London Math.
Soc. (2) \textbf{54} (1952), 147--167.

\bibitem{Wang22} L.~Wang, \emph{Identities on Zagier's rank two examples for Nahm's problem},
arXiv:2210.10748 (2022).

\bibitem{Warnaar50} S.~O. Warnaar, \emph{50 years of Bailey's lemma}, in: Algebraic Combinatorics and
Applications (A.~Betten et al., eds.), Springer, Berlin, 2001, 333--347; arXiv:0910.2062.

\bibitem{Warnaar2026} S.~O. Warnaar, \emph{Affine Jacobi--Trudi identities and $q,t$-Rogers--Ramanujan
identities}, SIGMA \textbf{22} (2026), 062; arXiv:2511.17034.

\end{thebibliography}
\end{document}
